% Generated by records_reorg.py assemble (records-reorg/); edit the data and re-run, not this file.
\subsection{Covariance construction and MINERvA's conventions}\label{recsec:methods-covariance}

\paragraph{MINERvA's published-covariance description.}
% [records reorg] former heading: \subsection{MINERvA's published-covariance description (from App. A, sec:minerva)}
% From app_statmethods.tex:235 @143f4fc97; context a reader does not need to interpret a result
% Framing added in the records (not moved text).
\noindent\emph{(``Below'' points to the note's MAT paragraph of \S\ref{sec:minerva}, which states the CV-centered pair formula.)}

\paragraph{Ruterbories 2021, Section V.}
% [records reorg] former heading: \paragraph{Ruterbories 2021, Section V.}
 The paper evaluates three uncertainty categories (flux, detector response and the neutrino-interaction model) by ``re-extracting the cross section using modified simulations'', i.e.\ the universe method, and gives no covariance formula. It reports a flux uncertainty below $4\%$ across most of phase space as a result of the $\nu$-$e$-scattering constraint (Valencia, arXiv:1906.00111); a $1.4\%$ target-nucleon normalization; an in-situ test-beam hadron response with per-particle GEANT inelastic variations; and a $1\%$ muon energy scale constrained by the low-recoil fit of Bashyal (JINST 16 P08068). It states that \emph{``the fit causes the flux and muon energy scale uncertainties to be correlated and that correlation is propagated to the final result''}; this is the only correlation between bands that the paper specifies. The only covariance in its equations is $V$ in the model-test $\chi^{2}$ (Eq.~1, ``the measurement covariance matrix''), with $205$ degrees of freedom (Table I); the paper also gives a log-normal $\chi^{2}$ variant, motivated by the highly correlated flux normalization (Peelle's Pertinent Puzzle). The paper itself gives no per-band covariance construction, no $\pm1\sigma$ pair formula and no band-independence or block-sum statement. MAT's code supplies the per-band construction and the independent band sum, but no separate pair formula (below). Nor does it state the rank, condition number and regularization of $V$, or any bootstrap or coverage validation (it uses D'Agostini iterative unfolding).

\paragraph{Upstream references of the band amplitudes.}
% [records reorg] former heading: \subsection{Upstream references of the band amplitudes (from App. A, sec:minerva)}
% From app_statmethods.tex:243 @143f4fc97; reference list

\paragraph{Upstream references that set each band's $\pm1\sigma$ amplitude.}
% [records reorg] former heading: \paragraph{Upstream references that set each band's $\pm1\sigma$ amplitude.}
\begin{itemize}[nosep]
  \item Flux / PPFX (100 universes), Aliaga, PRD 94 092005 (2016), arXiv:1607.00704;
  \item $\nu$-$e$ flux constraint ($<4\%$), Valencia, PRD 100 092001 (2019), arXiv:1906.00111;
  \item muon-E scale $\leftrightarrow$ flux fit, Bashyal, JINST 16 (2021) P08068;
  \item test-beam / GEANT hadron response, Aliaga, NIM A 789 28 (2015), arXiv:1501.06431;
  \item GENIE reweighting parameters, Andreopoulos, NIM A 614 87 (2010), and the GENIE manual arXiv:1510.05494;
  \item $1.4\%$ normalization (detector mass), Aliaga, NIM A 743 130 (2014), arXiv:1305.5199.
\end{itemize}

\paragraph{Multinomial to independent Poisson bootstrap limit.}
% [records reorg] former heading: \subsection{Multinomial to independent Poisson bootstrap limit (from App. A, app:poisson)}
% From app_statmethods.tex:126 @143f4fc97; derivation
Drawing $N$ events with replacement from an $N$-event sample gives counts $(n_{1},\dots,n_{N}) \sim \mathrm{Multinomial}(N,\mathbf{p})$ with uniform $p_{i}=1/N$ and $\sum_{i}n_{i}=N$. Each marginal $n_{i}\sim\mathrm{Binomial}(N,1/N)$ has mean $1$, variance $(N-1)/N \to 1$, and converges to $\mathrm{Poisson}(1)$ as $N$ grows (many trials, small success probability, fixed expected count). The $\sum n_{i}=N$ constraint induces $\mathcal{O}(1/N)$ correlations that vanish at large $N$. The multinomial is an independent Poisson process conditioned on $N$ counts. Removing that constraint gives independent Poisson(1) variables (Poissonization). One Poisson(1) factor is drawn per event. For $N \gtrsim 10^{6}$, our regime, the relative total-count error is $\mathcal{O}(1/\sqrt{N})$, far smaller than the per-bin Poisson variance.

\paragraph{Scope of the statistical appendix.}
% [records reorg] former heading: \subsection{From App. A, appstatmethods.tex: the paragraph beginning ``This appendix gives the statistical methods''}
% From app_statmethods.tex:5 @d55f7aad6; qualifier detail or procedure (DECISION-20261120a, second pass)
This appendix gives the statistical methods of this note's results, centered on the 2D comparison: covariance inversion; the training-seed, bootstrap and universe estimators, with the Poisson limit and the $\pm1\sigma$ pair covariance; their MAT-conformant construction and combination; the $\chi^{2}$, pull and shape-only conventions; the rank and conditioning of the covariances; the anatomy of the tension in the 2D comparison; the closure procedures; and a summary of the 2D comparison numbers. It covers analysis-specific conventions, not general statistics. The $\chi^{2}$ convention used throughout is Eq.~\eqref{eq:chi2} in \S\ref{sec:chisq}.

\paragraph{Coherent Poisson bootstrap.}
% [records reorg] former heading: \subsection{From App. A, appstatmethods.tex: the paragraph beginning ``A coherent bootstrap carries each data''}
% From app_statmethods.tex:103 @d55f7aad6; qualifier detail or procedure (DECISION-20261120a, second pass)
In the paragraph as the note printed it, this text followed: ``A coherent bootstrap carries each data factor through step-1 classifier training and final cross-section construction.''\par
Each MC factor is shared by the reco-tree and truth-tree copies, preserving migration consistency. The three sub-RNGs (data, signal MC, background MC; Table~\ref{tab:variance-sources}) make the sources uncorrelated by construction, and pinning the GBDT seed isolates the bootstrap from the ML noise of \S\ref{sec:seedscan}. A Poisson(1) factor can be zero. Dropping those events from each replica mimics a resample with replacement. Poisson(1) is the large-$N$ limit of the multinomial counts (\S\ref{app:poisson}).

\paragraph{Bootstrap replica count.}
% [records reorg] former heading: \subsection{From App. A, appstatmethods.tex: the paragraph beginning ``For these replicas (205 bins, MEFHC''}
% From app_statmethods.tex:117 @d55f7aad6; qualifier detail or procedure (DECISION-20261120a, second pass)
% [lead-in, Option B] In the paragraph as the note printed it, this text followed: ``For these replicas (205 bins, ME~FHC LightGBM, seed~42) the relative standard deviation of the total cross section is $\boxed{0.061\%}$, the per-bin median relative spread $\boxed{0.676\%}$, and the per-bin $p_{84}$ $1.602\%$.''\par
With $300$ replicas the per-bin $\sigma$ has a relative standard error of about $4.1\%$ by the Gaussian formula $1/\sqrt{2(N-1)}$. $N=300$ is chosen for full rank.

\paragraph{Bands in the universe sweep.}
% [records reorg] former heading: \subsection{From App. A, appstatmethods.tex: the paragraph beginning ``The universe sweep used by the''}
% From app_statmethods.tex:130 @d55f7aad6; qualifier detail or procedure (DECISION-20261120a, second pass)
% [lead-in, Option B] In the paragraph as the note printed it, this text followed: ``The universe sweep used by the quoted budget has 187 universes across 44 bands, dominated by the $100$-universe PPFX flux band (hadron-production sampling).''\par
The quoted band sum drops one band (Rvp1pi) and holds $43$ bands plus the normalization.
% [records reorg] former heading: \subsection{From App. A, appstatmethods.tex: the paragraph beginning ``The universe sweep used by the''}
% From app_statmethods.tex:130 @d55f7aad6; qualifier detail or procedure (DECISION-20261120a, second pass)
In the paragraph as the note printed it, this text followed: ``Almost all other bands are $\pm1\sigma$ pairs ($N_{u}=2$).''\par
They include the MINOS and lateral MINERvA muon energy scales, the MINOS muon-match efficiency, the GENIE parameters \texttt{MaRES}, \texttt{MvRES} and \texttt{MaCCQE}, and the lateral GEANT FSI responses for neutrons, pions and protons. Where selection migration matters, the five kinematic bands (the two beam-angle components, the muon resolution and the two muon energy scales) require re-selected event loops, run with one universe active at a time (weight-only bands: \S\ref{sec:uq-higherdim}). The quoted universe block takes these five bands from re-unfolds of re-selected universes.

\paragraph{Two-universe limit of the MAT covariance.}
% [records reorg] former heading: \subsection{From App. A, appstatmethods.tex: the paragraph beginning ``MINERvA's MAT framework (textttPlotUtils/allowbreak MnvVertErrorBand::allowbreak CalcCovMx,''}
% From app_statmethods.tex:146 @d55f7aad6; qualifier detail or procedure (DECISION-20261120a, second pass)
% [lead-in, Option B] In the paragraph as the note printed it, this text followed: ``uniformly for all $N_{u}\ge 2$, with no special $\pm1\sigma$ pair formula.''\par
At $N_{u}=2$ it reduces to rank-1 $\tfrac14(\mathbf{X}_{+}-\mathbf{X}_{-})(\mathbf{X}_{+}-\mathbf{X}_{-})\T$. Mean-centering removes a common offset from the central value. The half-difference remains and need not be small. For a one-sided or asymmetric pair, CV-centering adds that offset as a second direction (\S\ref{app:pair}).

\paragraph{Corrected systematic construction.}
% [records reorg] former heading: \subsection{From App. A, appstatmethods.tex: the paragraph beginning ``Equationeqrefeq:systblocksum omits a larger correction. The''}
% From app_statmethods.tex:159 @d55f7aad6; qualifier detail or procedure (DECISION-20261120a, second pass)
% [lead-in, Option B] In the paragraph as the note printed it, this text followed: ``A unified throw measures it through $N$ re-unfolds, simultaneously perturbing every reweightable band and the flux multisim per event ($w_{\CV}\prod_B\rho_B^{g_B}$, $g_B\sim\mathcal N(0,1)$), then re-extracting the cross section.''\par
The corrected construction follows. Its use in the five-axis covariance of the \resFiveAxis{} is in \S\ref{sec:eavailw-adopted} (Eq.~\eqref{eq:zinflation}).
In the note, the corrected construction that follows is the unified throw, described in the next paragraph of App.~\ref{app:statmethods}.\par
% [records reorg] former heading: \subsection{From App. A, appstatmethods.tex: the paragraph beginning ``Grouped by category as in Fig.reffig:uqbands,''}
% From app_statmethods.tex:199 @d55f7aad6; qualifier detail or procedure (DECISION-20261120a, second pass)
% [lead-in, Option B] In the paragraph as the note printed it, this text followed: ``Grouped by category as in Fig.~\ref{fig:uqbands}, which follows Figs.~6--7 of Ref.~\cite{MINERvA:2021owq}, the medians are: Muon reconstruction $3.437\%$, Models $1.428\%$, Flux $3.723\%$, Hadronic response $0.752\%$, Statistical $0.676\%$, Normalization $1.400\%$, Training seed $0.048\%$, Run period $2.087\%$.''\par
(MINOS$+$MINERvA) and the MINOS-efficiency band dominate; the first is the band that the Ruterbories paper identifies as flux-correlated (\S\ref{sec:rank}).

\paragraph{Mean- versus CV-centered pair estimators.}
% [records reorg] former heading: \subsection{From App. A, appstatmethods.tex: the paragraph beginning ``For two universes u=0,1 at pm''}
% From app_statmethods.tex:227 @d55f7aad6; qualifier detail or procedure (DECISION-20261120a, second pass)
The estimators agree only for CV-symmetric pairs ($\mathbf{m}=0$), both reducing to $\mathbf{a}\mathbf{a}\T$. For asymmetric or one-sided pairs, CV-centering adds the rank-1 $\mathbf{m}\mathbf{m}\T$ common-offset direction. Every per-band block $C_{b}$ here is mean-centered to match MAT (\S\ref{sec:minerva}); this matters for BeamAngle, MuonResolution and the one-sided NC band.

\paragraph{MAT covariance convention.}
% [records reorg] former heading: \subsection{From App. A, appstatmethods.tex: the paragraph beginning ``MAT's C++ framework defines the covariance''}
% From app_statmethods.tex:233 @d55f7aad6; qualifier detail or procedure (DECISION-20261120a, second pass)
MAT's C++ framework defines the covariance construction that the published paper leaves unspecified. It is the mean-centered $1/N$ per-band sample covariance and element-wise band sum used here.

\paragraph{MAT implementation and normalization.}
% [records reorg] former heading: \subsection{From App. A, appstatmethods.tex: the paragraph beginning ``textttMnvVertErrorBand::CalcCovMx implements Eq.eqrefeq:matcov, referenced to the''}
% From app_statmethods.tex:236 @d55f7aad6; qualifier detail or procedure (DECISION-20261120a, second pass)
\texttt{MnvVertErrorBand::CalcCovMx} implements Eq.~\eqref{eq:matcov}, referenced to the universe mean $\langle x\rangle_{u}$ rather than $X_{\CV}$, normalized by biased $1/N$ rather than Bessel's $1/(N-1)$ (a $1\%$ offset for PPFX, $N=100$), with no special case for $N=2$. The conventional CV-centered pair formula (``MINERvA-101 $\pm1\sigma$ pair''), $\tfrac12(\boldsymbol{\Delta}_{+}\boldsymbol{\Delta}_{+}\T + \boldsymbol{\Delta}_{-}\boldsymbol{\Delta}_{-}\T)$, agrees with it only for CV-symmetric shifts. The quoted universe block, filled into a \texttt{PlotUtils::MnvH1D} with MAT's universe weights on Flux, agrees with \texttt{MnvH1D::GetTotalErrorMatrix} to $1.3\times10^{-16}$ of its largest element. The quoted Flux band carries MAT's $\nu$--$e$ and inverse-muon-decay constraint weights.

\paragraph{Band sum and rank of the total covariance.}
% [records reorg] former heading: \subsection{From App. A, appstatmethods.tex: the paragraph beginning ``textttMnvH1D::GetTotalErrorMatrix sums bands and the statistical''}
% From app_statmethods.tex:240 @d55f7aad6; qualifier detail or procedure (DECISION-20261120a, second pass)
% [lead-in, Option B] In the paragraph as the note printed it, this text followed: ``\texttt{MnvH1D::GetTotalErrorMatrix} sums bands and the statistical matrix element-wise, assuming band independence.''\par
Cross-band correlations exist only if encoded in a combined band, such as the Bashyal joint fit handled outside this loop. MAT applies no ridge, shrinkage, SVD truncation or rank-deficiency handling to the published total covariance. The statistical block makes $V$ full-rank. Its systematic-only rank is $120$, below the $138$ of our universe block (eigenvalues above $10^{-12}\lambda_{\max}$; \S\ref{sec:rank}).

\paragraph{Variance sources of the blocks.}
% [records reorg] former heading: \subsection{From App. A, appstatmethods.tex: the paragraph beginning ``The bootstrap varies events through distinct''}
% From app_statmethods.tex:245 @d55f7aad6; qualifier detail or procedure (DECISION-20261120a, second pass)
The bootstrap varies events through distinct data, signal-MC and background-MC sub-RNGs; universes vary model weights at a pinned GBDT seed; the training-seed block varies the GBDT seed over five settings, each on the same $20$ bootstrap replicas (at seed~42, replicas of the statistical block); the run-period block compares separate unfolds of three run periods. The bootstrap and the universes use the central value's LightGBM seed ($42$), and the training-seed settings include it.

\paragraph{Combined median uncertainties.}
% [records reorg] former heading: \subsection{From App. A, appstatmethods.tex: the paragraph beginning ``The combined sqrtTrmathbfC of Tablereftab:uq is''}
% From app_statmethods.tex:262 @d55f7aad6; qualifier detail or procedure (DECISION-20261120a, second pass)
% [lead-in, Option B] In the paragraph as the note printed it, this text followed: ``The relative entries divide by the quoted central value, the LightGBM seed-42 central value of the universe sweep.''\par
Dividing the combined covariance by the exact-split cross-check instead gives a combined median of \SI{6.468}{\percent}. The universe, statistical and training-seed blocks are LightGBM ensembles, the first two at the central value's training seed and the third over five settings that include it (\S\ref{sec:method-config}). The universe and statistical blocks together give $5.976\%$, again evaluated on their sum rather than combined from the printed medians.
% [records reorg] former heading: \subsection{From App. A, appstatmethods.tex: the paragraph beginning ``The data/MC split bootstrap was run''}
% From app_statmethods.tex:431 @d55f7aad6; qualifier detail or procedure (DECISION-20261120a, second pass)
% [lead-in, Option B] In the paragraph as the note printed it, this text followed: ``For the quoted statistical block, $\sqrt{\Tr\mathbf{C}^{\boot}}/\sqrt{\Tr\,\texttt{StatOnly}}=0.407$ and the per-bin median ratio of the two $\sigma$ is $0.425$.''\par
The fixed-truth test (\S\ref{sec:coverage}, \S\ref{sec:uq-boot}) does not decide between the explanations above.

\paragraph{Variance versus bias.}
% [records reorg] former heading: \subsection{From App. A, appstatmethods.tex: the paragraph beginning ``The variance methods of Srefsec:variance quantify''}
% From app_statmethods.tex:459 @d55f7aad6; qualifier detail or procedure (DECISION-20261120a, second pass)
The variance methods of \S\ref{sec:variance} quantify the scatter of the unfolded result. Closure tests measure its bias, a consistent offset from a known injected truth, which the variance methods cannot show.

\paragraph{Acceptance leakage bound.}
% [records reorg] former heading: \subsection{From App. A, appstatmethods.tex: the paragraph beginning ``OmniFold sees reco and truth (pT,pPar)''}
% From app_statmethods.tex:475 @d55f7aad6; qualifier detail or procedure (DECISION-20261120a, second pass)
% [lead-in, Option B] In the paragraph as the note printed it, this text followed: ``For median per-bin signed residual $\bar r(A)$, a linear fit gives $\bar r(A) \approx \kappa_{\text{leak}} A$, with $\kappa_{\text{leak}} = \partial \bar r/\partial A$, the fraction of hidden $\dpT$ bias entering the visible answer.''\par
The $3.0\%$ acceptance bound is looser than the truth-reweight bound because the weak $\dpT$--$\pT$ correlation prevents complete removal of the leakage.

\paragraph{Normality of toy residuals.}
% [records reorg] former heading: \subsection{From App. A, appstatmethods.tex: the paragraph beginning ``In reported bin i with bar''}
% From app_statmethods.tex:520 @d55f7aad6; qualifier detail or procedure (DECISION-20261120a, second pass)
% [lead-in, Option B] In the paragraph as the note printed it, this text followed: ``The fraction with $|r|\le1$ is \SI{\pullFrac}{\percent} (Gaussian reference $68.27\%$; $\Delta = +0.44\%$; median $68.50\%$), and $\langle|r|\rangle=\pullMeanAbs$ against $\sqrt{2/\pi}=\pullGauss$ (within $0.5\%$).''\par
The signed residual mean is $+0.006 \pm 0.082\,\sigma$, consistent with zero. The within-bin fraction is above \SI{65}{\percent} in \SI{97.6}{\percent} of bins (only $5$ of $205$ fall below). These show that the standardized toy residuals are approximately normal.
% [records reorg] former heading: \subsection{From App. A, appstatmethods.tex: the paragraph beginning ``Because sigmai is estimated from the''}
% From app_statmethods.tex:522 @d55f7aad6; qualifier detail or procedure (DECISION-20261120a, second pass)
% [lead-in, Option B] In the paragraph as the note printed it, this text followed: ``Because $\sigma_i$ is estimated from the same toys that are evaluated, these figures diagnose standardized pulls and Gaussianity only and are not a coverage test of the bootstrap band.''\par
No coverage number is quoted from this ensemble, and the 2D coverage statements are the fixed-truth test of \S\ref{sec:coverage} and the held-out-tune test of \S\ref{sec:coverage-tunes}.
% carried with the text from app_statmethods.tex:523 @d55f7aad6
% Search: coverage-toys; coverage_summary; same-ensemble pull @ ddf16f10

\subsection{Numerical conventions, inversion and rank}\label{recsec:methods-rank}

\paragraph{Regularization scans on an earlier covariance.}
% [records reorg] former heading: \subsection{Regularization scans on an earlier covariance (from App. A, sec:rank)}
% From app_statmethods.tex:465 @143f4fc97; table of an earlier covariance and the exact-split cross-check
% NR: no committed producer of this scan (rebuild REPORT §6); kept as a labelled historical diagnostic.
% J3/J4 provisional (DECISION-20261018, packet C7: kept, labelled historical).
\begin{center}
\resizebox{\textwidth}{!}{%
\begin{tabular}{@{}lll@{}}
\toprule
Method & Setting & Result \\
\midrule
Eigenvalue truncation & $\lambda/\lambda_{\max} > \{10^{-3},10^{-4},10^{-5}\}$ &
  rank $\{19,49,82\}$; $\chi^{2}/205 = \{0.13,1.12,3.08\}$ \\
Ridge $\mathbf{C}+\varepsilon\lambda_{\max}\mathbf{I}$ & $\varepsilon=\{10^{-3},10^{-4},10^{-5}\}$ &
  $\chi^{2}/205 = \{0.35,1.38,4.30\}$ \\
Diagonal-target shrinkage & $\lambda=0.05$ &
  rank $178\to204$, cond $1.3\times10^{11}$, $\chi^{2}/205 = 3.29$ \\
\bottomrule
\end{tabular}}
\end{center}

\paragraph{Rank of sample covariances.}
% [records reorg] former heading: \subsection{From App. A, appstatmethods.tex: the paragraph beginning ``A sample covariance from Nle n''}
% From app_statmethods.tex:20 @d55f7aad6; qualifier detail or procedure (DECISION-20261120a, second pass)
A sample covariance from $N\le n$ replicas in $n$ bins has rank at most $N-1$.
% [records reorg] former heading: \subsection{From App. A, appstatmethods.tex: the paragraph beginning ``A sample covariance from Nle n''}
% From app_statmethods.tex:20 @d55f7aad6; qualifier detail or procedure (DECISION-20261120a, second pass)
% [lead-in, Option B] In the paragraph as the note printed it, this text followed: ``The 300-replica bootstrap ($N=300 > 205$) has more replicas than bins, which permits full rank and does not establish it.''\par
The rank of the systematic block is limited at any universe count (\S\ref{sec:rank}).

\paragraph{Global scale of the bootstrap block.}
% [records reorg] former heading: \subsection{From App. A, appstatmethods.tex: the paragraph beginning ``Even at full rank (205) a''}
% From app_statmethods.tex:421 @d55f7aad6; qualifier detail or procedure (DECISION-20261120a, second pass)
The bootstrap block is smaller in \emph{global scale}: its square-root trace is $0.407$ of the paper's ($\sqrt{\Tr\mathbf{C}^{\boot}}=1.886\times10^{-40}$ vs $\sqrt{\Tr\,\texttt{StatOnlyCov}}=4.639\times10^{-40}$).
% [records reorg] former heading: \subsection{From App. A, appstatmethods.tex: the paragraph beginning ``Even at full rank (205) a''}
% From app_statmethods.tex:421 @d55f7aad6; qualifier detail or procedure (DECISION-20261120a, second pass)
% [lead-in, Option B] In the paragraph as the note printed it, this text followed: ``The direct inverse of $\mathbf{C}^{\syst}+\mathbf{C}^{\boot}$ gives $\chi^{2}/\mathrm{ndf}=22.38$ ($21.96$ for all four blocks) for the LightGBM central value, against the paper covariance's $\chiPaper$ with its larger stat floor, so a usable $\chi^{2}_{\ours}$ depends on the \emph{statistical} block, not on more systematic bands.''\par
The gap is \emph{not} a missing MC-statistical term: our bootstrap already draws independent Poisson($1$) multipliers on data, signal MC and background MC per replica (\S\ref{sec:bootstrap}; Practical Guide 2507.09582 \S5.1). Three candidates remain. The first is a real difference in statistical efficiency between the unbinned reweighting and the published binned D'Agostini unfold. The second is a difference in what the paper's \texttt{StatOnlyCov} encompasses. The third is a regularization effect in our own estimator that suppresses dispersion without improving information. The split-bootstrap diagnostic used on the earlier rebuilt bootstrap has not been repeated for the quoted statistical block.

\paragraph{Seed pinning and cross-correlations.}
% [records reorg] former heading: \subsection{From App. A, appstatmethods.tex: the paragraph beginning ``Eq.eqrefeq:blocksum requires no cross-correlations between ML,''}
% From app_statmethods.tex:499 @d55f7aad6; qualifier detail or procedure (DECISION-20261120a, second pass)
% [lead-in, Option B] In the paragraph as the note printed it, this text followed: ``\textbf{Independence of error sources.} Eq.~\eqref{eq:blocksum} requires no cross-correlations between ML, statistical, systematic and run-period blocks.''\par
The bootstrap pins the GBDT seed, and universes pin the seed too. A common seed correlates the training noise and does not cancel it (\S\ref{sec:uq-2d-scope}); the training-seed block averages the same $20$ bootstrap replicas at each seed setting and includes the bootstrap\,$\times$\,seed interaction at its $1/20$ share. Among the first three, cross-terms are higher-order.

\subsection{Comparison statistics, closure and coverage}\label{recsec:methods-comparison}

\paragraph{Log-normal chi2 definitions compared on an earlier covariance.}
% [records reorg] former heading: \subsection{Log-normal chi2 definitions compared on an earlier covariance (from App. A, sec:chisq)}
% From app_statmethods.tex:335 @143f4fc97; history of an earlier covariance
% Framing added in the records (not moved text); the J3/J4 comment that the moved "comment below" names stays in app_statmethods.tex.
\noindent\emph{(``The definition above'' is the note's log-normal definition in \S\ref{sec:chisq}.)}
 An independent recomputation compared five definitions on an earlier covariance, before the refined estimator was adopted. Only the definition above reproduced the value then quoted. The other four gave $1.197$ (our covariance scaled by our own values), $0.958$ (both scaled by ours), $0.963$ (scaled by the geometric mean) and $1.015$ ($\ln(1+\mathbf{C}/(x_{i}x_{j}))$ in place of the ratio). The quoted $\chiCombined$ and $\chiCombinedLog$ have no independent recomputation: they are the production run's own outputs. Its checks confirm that the covariance's reported cells are the paper's and that the plotted projections equal the budget's, and neither recomputes these values. Adding the OmniFold error budget lowers the summed-covariance distance from $\chiPaper$ to $\chiCombined$ per bin. % G2R no-recompute sentence (the director, 2026-11-15): RE-2-shapes run 2 has no independent chi2 recompute (cov2d, 2026-11-15); its check_rebuild (g2r/logs/check_rebuild.log, 5/5) checks the chi2 files' cell identity (chi2_Q_std.txt: 'cell identity OK: ... hReportedCells equals the paper StatOnly set') and fig6/7 against budget_Q. The five-definition recompute is the Stage F rebuild's, on its Q (comment below).

\paragraph{Edge-bin sensitivity table.}
% [records reorg] former heading: \subsection{Edge-bin sensitivity table (from App. A, sec:chisq)}
% From app_statmethods.tex:369 @143f4fc97; table of a robustness check of the exact-split cross-check
% (\texttt{compare\_to\_paper\_fullcov.py})
\begin{center}
\begin{tabular}{@{}rrrrr@{}}
\toprule
$\pPar \ge$ (GeV/c) & $N_{\text{bins}}$ & $\chi^{2}$ & $\chi^{2}/\mathrm{ndf}$ & median ratio \\
\midrule
1.5 & 205 & 750.49 & 3.661 & 1.0064 \\
2.0 & 196 & 685.36 & 3.497 & 1.0059 \\
2.5 & 186 & 545.87 & 2.935 & 1.0059 \\
3.0 & 175 & 474.84 & 2.713 & 1.0059 \\
3.5 & 163 & 452.95 & 2.779 & 1.0073 \\
4.0 & 151 & 428.72 & 2.839 & 1.0068 \\
\bottomrule
\end{tabular}
\end{center}

\paragraph{Sources of the 2D comparison numbers.}
% [records reorg] former heading: \subsection{Sources of the 2D comparison numbers (from App. A, sec:assemble)}
% From app_statmethods.tex:536 @143f4fc97; an index of numbers stated in their own sections

\paragraph{Sources of the 2D comparison numbers.}
% [records reorg] former heading: \subsubsection*{Sources of the 2D comparison numbers}
\begin{center}
\begin{tabularx}{\textwidth}{@{}>{\RaggedRight\arraybackslash}p{0.27\textwidth}l>{\RaggedRight\arraybackslash}X@{}}
\toprule
Number & Section & Procedure \\
\midrule
$\chi^{2}/\mathrm{ndf} = \chiCombined$ & \S\ref{sec:chisq} &
  Eq.~\eqref{eq:chi2} with $\mathbf{C}^{\paper}+\mathbf{C}^{\syst}+\mathbf{C}^{\stat}+\mathbf{C}^{\ML}+\mathbf{C}^{\mathrm{PL}}$ \\
$\mathbf{C}^{\syst}$ & \S\ref{sec:universes} &
  Universe sweep ($187$ universes; $43$ bands in the band sum), MAT-conformant Eq.~\eqref{eq:matcov} $+$ norm \\
$\mathbf{C}^{\stat}$ & \S\ref{sec:bootstrap} &
  Poisson bootstrap, $300$ replicas, seed $42$, three streams, Eq.~\eqref{eq:bootcov} \\
$\mathbf{C}^{\ML}$ & \S\ref{sec:seedscan} & LightGBM training-seed settings, $5\times20$ replicas \\
$\mathbf{C}^{\mathrm{PL}}$ & \S\ref{sec:uq-combine} & Run-period block, three run periods \\
$\mathbf{C}^{\paper}$ & \S\ref{sec:minerva} & Paper's published \texttt{TotalCovariance} \\
Pull standard deviation $= 0.389$ & \S\ref{sec:pull} & % G2R: RE-2-shapes run 2 (C_new), g2r-integration/ledger.tsv g2r-audit-pull-rms-b; a source comment citing RE-2 part C for this value gives the earlier value
  Per-bin standardized residual on combined $\mathbf{C}$ \\
Combined median rel.\ $= 6.463\%$ & \S\ref{sec:combine} & % G2R: RE-2-shapes run 2 (C_new), g2r-integration/ledger.tsv g2r-audit-q-median; a source comment citing RE-2 part C for this value gives the earlier value
  $\sqrt{\diag\mathbf{C}^{\tot}}/\mathbf{X}^{\CV}$, $\mathbf{X}^{\CV}$ the LightGBM seed-42 central value \\
Fixed-truth $C_2=\covTwoDCTwo$ (nominal $0.954$) & \S\ref{sec:coverage} &
  Earlier statistical band vs.\ fixed truth, $N_{\text{toy}}=\covTwoDNtoy$: fails \\
% J3/J4 provisional (DECISION-20261018, packet W3: D2): the two coverage rows score earlier bands, not Q's statistical block (not recomputed, NR).
Same-ensemble $|r|\leq1$: $68.71\%$ & \S\ref{sec:pull-diagnostic} &
  Pull diagnostic, $N_{\text{toy}}=200$ \\
\bottomrule
\end{tabularx}
\end{center}

\paragraph{Same-ensemble pull diagnostic: the toy ensemble.}
% [records reorg] former heading: \subsection{Same-ensemble pull diagnostic: the toy ensemble (from App. A, sec:pull-diagnostic)}
% From app_statmethods.tex:514 @e57f0a0a8; procedure
 Each toy has its own bootstrap seed $k$ and per-event Poisson(1) draws, as in production, with seeds chosen not to collide with the production replicas. Toy seeds are independent. Bins within a toy are correlated. The MC bootstrap multiplies the truth weight. It fluctuates both the unfolded result and the stored truth. The stored truth varies in all 205 reported bins, and its all-toy mean $\bar X^{\text{truth}}$ is therefore not an independent fixed truth for a split-sample coverage claim. It only centers a same-ensemble diagnostic.

\paragraph{Bins suppressed by the efficiency cut.}
% [records reorg] former heading: \subsection{From App. A, appstatmethods.tex: the paragraph beginning ``The paper reports mathrmd2sigma/(mathrmdpT,mathrmdpPar) on 14''}
% From app_statmethods.tex:282 @d55f7aad6; qualifier detail or procedure (DECISION-20261120a, second pass)
% [lead-in, Option B] In the paragraph as the note printed it, this text followed: ``Of these, $205$ have $x^{\paper}>0$ in \texttt{data\_result\_ptpl\_2D\_minerva\_inclusive\_6GeV.txt}.''\par
The paper suppresses the remaining $19$ through an efficiency cut. Every $\chi^{2}$ and pull excludes them on both sides.

\paragraph{Strict-interior mask.}
% [records reorg] former heading: \subsection{From App. A, appstatmethods.tex: the paragraph beginning ``A secondary strict-interior mask, used only''}
% From app_statmethods.tex:284 @d55f7aad6; qualifier detail or procedure (DECISION-20261120a, second pass)
A secondary strict-interior mask, used only as a sensitivity check, keeps cells with $\pT^{\text{high}} / \pPar^{\text{low}} \le \tan 20^{\circ}$, removing the $20$ reported cells that straddle the $\theta_{\mu}<20^{\circ}$ kinematic cut. The $185$-bin check tests whether the main tension comes from the boundary.
% carried with the text from app_statmethods.tex:285 @d55f7aad6
% Source: pre-freeze 2D_OMNIFOLD_RUN_LOG_ARCHIVE.md:50-59 (Phase 4, unreported bins), :716-726.

\paragraph{Computation of the comparison $\chi^{2}$.}
% [records reorg] former heading: \subsection{From App. A, appstatmethods.tex: the paragraph beginning ``With boldsymbolDeltaequivmathbfxours-mathbfxpaper on the 205-bin set''}
% From app_statmethods.tex:305 @d55f7aad6; qualifier detail or procedure (DECISION-20261120a, second pass)
The comparison applies the pseudo-inverse of \S\ref{sec:invert} to the $205\times205$ block and computes the quadratic form. Options add OmniFold covariances via Eq.~\eqref{eq:blocksum} or subtract the paper's stat block to avoid counting it alongside the bootstrap (\S\ref{sec:rank}).

\paragraph{Truncation scan and ndf choice.}
% [records reorg] former heading: \subsection{From App. A, appstatmethods.tex: the paragraph beginning ``Setting mathrmndf=ntextreported requires nearly that many''}
% From app_statmethods.tex:309 @d55f7aad6; qualifier detail or procedure (DECISION-20261120a, second pass)
The choice rests on the $205$-bin 2D truncation scan rather than being assumed: the scan rises smoothly with the retained rank $r$, $0.39\,(r{=}50)\to1.04\,(100)\to1.76\,(180)\to2.27\,(205)$ (exact-split cross-check: $0.69\to2.35\to3.30\to3.66$), with no plateau or late jump, indicating an effective rank comparable to the dimension. This evidence is specific to 2D and does not transfer to N-D covariances (\S\ref{sec:nd-declarations}).

\paragraph{Missing cross-covariance.}
% [records reorg] former heading: \subsection{From App. A, appstatmethods.tex: the paragraph beginning ``emph(i) Missing cross-covariance prevents calibration. For''}
% From app_statmethods.tex:338 @d55f7aad6; qualifier detail or procedure (DECISION-20261120a, second pass)
% [lead-in, Option B] In the paragraph as the note printed it, this text followed: ``and the published IBU-era and OmniFold results share MINERvA data and flux, detector and GENIE systematics.''\par
Without $C_{xy}$, the summed $C_{x}+C_{y}$ double-counts shared systematics and omits their correlation. Its net error has unknown size and sign.

\paragraph{Expected spread of $\chi^{2}/\mathrm{ndf}$.}
% [records reorg] former heading: \subsection{From App. A, appstatmethods.tex: the paragraph beginning ``emph(ii) Near-one realizations are not evidence''}
% From app_statmethods.tex:340 @d55f7aad6; qualifier detail or procedure (DECISION-20261120a, second pass)
Under calibrated assumptions, $\chi^{2}/\mathrm{ndf}$ has expectation $1$ and spread of order $\sqrt{2/\mathrm{ndf}}$, about $0.10$ for $\mathrm{ndf}=205$.

\paragraph{Shared estimator and scope of the scalar.}
% [records reorg] former heading: \subsection{From App. A, appstatmethods.tex: the paragraph beginning ``emph(iii) It does not address coverage''}
% From app_statmethods.tex:342 @d55f7aad6; qualifier detail or procedure (DECISION-20261120a, second pass)
% [lead-in, Option B] In the paragraph as the note printed it, this text followed: ``\emph{(iii) It does not address coverage} (\S\ref{sec:uq-2d-scope}, \S\ref{sec:coverage}).''\par
The covariance and the central value share the LightGBM estimator and, in the universe and statistical blocks, its training seed (\S\ref{sec:method-config}). Agreement claims use per-bin standardized differences and ratios (\S\ref{sec:pull}); this scalar is only a distance and ratio to the paper. The log-normal value $\chiCombinedLog$ depends on the definition given with the worked example.

\paragraph{Statistical floor in the combined covariance.}
% [records reorg] former heading: \subsection{From App. A, appstatmethods.tex: the paragraph beginning ``The quoted chiCombined inverts the emphsum''}
% From app_statmethods.tex:349 @d55f7aad6; qualifier detail or procedure (DECISION-20261120a, second pass)
The quoted $\chiCombined$ inverts the \emph{sum} of the paper's covariance and ours, not our blocks alone. That sum contains the paper's \texttt{StatOnlyCov}, its published statistical-only covariance, which has rank $205$ and is near-diagonal (\S\ref{sec:rank}), so it places a variance floor under \emph{every} eigendirection of the combined matrix. It is the larger of the two statistical contributions. Our statistical block is $0.407$ of it in \emph{square-root} trace, and the square of that factor in \emph{trace} (operands in \S\ref{sec:rank}). Each ratio states which power of the trace it uses, because the two ratios describe the same comparison and differ by more than a factor of two. An external, published block, rather than any ensemble of ours, therefore sets the small-eigenvalue directions on which a finite-$N$ debiasing factor would act. This bounds the quoted number's exposure to declaration~(v) of \S\ref{sec:nd-declarations}.

\paragraph{Trace share and inverse dominance.}
% [records reorg] former heading: \subsection{From App. A, appstatmethods.tex: the paragraph beginning ``Trace share alone cannot justify this''}
% From app_statmethods.tex:354 @d55f7aad6; qualifier detail or procedure (DECISION-20261120a, second pass)
Trace share alone cannot justify this protection: trace weights eigenvalues by $\lambda$, precision by $1/\lambda$, so even a negligible-trace block can dominate an inverse. Without the external floor, our finite-$N$ blocks fill $\mathbf{C}^{\syst}$'s null space (rank $138\to205$, eigenvalues above $10^{-12}\lambda_{\max}$) and, inverted directly, yield $\chi^{2}/\mathrm{ndf}=22.38$ (\S\ref{sec:rank}). Protection depends on the source of the smallest retained eigenvalues.
% [records reorg] former heading: \subsection{From App. A, appstatmethods.tex: the paragraph beginning ``The paper covariance alone gives mean''}
% From app_statmethods.tex:382 @d55f7aad6; qualifier detail or procedure (DECISION-20261120a, second pass)
% [lead-in, Option B] In the paragraph as the note printed it, this text followed: ``The mean is a small difference between two estimates of the same data. It is not a finding of no bias.''\par
The standard deviation is slightly larger before our uncertainties are added.

\paragraph{Shape-only comparison.}
% [records reorg] former heading: \subsection{From App. A, appstatmethods.tex: the paragraph beginning ``This tests whether the emphshape of''}
% From app_statmethods.tex:387 @d55f7aad6; qualifier detail or procedure (DECISION-20261120a, second pass)
This tests whether the \emph{shape} of the $(\pT,\pPar)$ cross section matches the paper independently of absolute normalization, which depends on the flux integral, target nucleon count, and POT (each can drift by $\mathcal{O}(\%)$); the shape is a pure unfolding diagnostic.

\paragraph{Band count and covariance rank.}
% [records reorg] former heading: \subsection{From App. A, appstatmethods.tex: the paragraph beginning ``A data-to-model chi2/mathrmndf with this analysis's''}
% From app_statmethods.tex:409 @d55f7aad6; qualifier detail or procedure (DECISION-20261120a, second pass)
% [lead-in, Option B] In the paragraph as the note printed it, this text followed: ``Every rank and condition number of our blocks quoted below uses this definition.''\par
Band count, not universe count, limits the rank: Eq.~\eqref{eq:matcov} gives each $\pm1\sigma$ pair rank $1$ and the $100$-universe PPFX band rank $\le 99$. PPFX is effectively rank-$1$ (\S\ref{sec:uq-universes}). The $43$ bands cannot span $205$ bins. The deficit does not come from missing bands.
% carried with the text from app_statmethods.tex:410 @d55f7aad6
% Source of rank and cond: budget_Q ranks (review NOTE 1). The ours-only script prints its own cond for the
% [records reorg] former heading: \subsection{From App. A, appstatmethods.tex: the paragraph beginning ``A data-to-model chi2/mathrmndf with this analysis's''}
% From app_statmethods.tex:409 @d55f7aad6; qualifier detail or procedure (DECISION-20261120a, second pass)
% [lead-in, Option B] In the paragraph as the note printed it, this text followed: ``The paper's systematic-only covariance ($\texttt{TotalCov}-\texttt{StatOnlyCov}$) has rank $120$, below our $138$.''\par
We use the same \texttt{GetStandardSystematics} and propagate at least as many independent systematic directions.
% carried with the text from app_statmethods.tex:410 @d55f7aad6
% Source of rank and cond: budget_Q ranks (review NOTE 1). The ours-only script prints its own cond for the
% [records reorg] former heading: \subsection{From App. A, appstatmethods.tex: the paragraph beginning ``The combined distance (Srefsec:chisq) adds mathbfCboot''}
% From app_statmethods.tex:442 @d55f7aad6; qualifier detail or procedure (DECISION-20261120a, second pass)
In the paragraph as the note printed it, this text followed: ``Subtracting that block uses $\texttt{TotalCov}-\texttt{StatOnlyCov}$ as baseline before our additions, changing $\chi^{2}/\mathrm{ndf}$ $\chiCombined\to \chiCombinedSubStat$.''\par
It over-corrects: our $\sqrt{\Tr\mathbf{C}^{\boot}}=1.886\times10^{-40}$ is $0.407$ of the paper's $\sqrt{\Tr\,\texttt{StatOnlyCov}}=4.639\times10^{-40}$, leaving no comparable replacement.

\paragraph{Truncated pseudo-inverse as collaboration practice.}
% [records reorg] former heading: \subsection{From App. A, appstatmethods.tex: the paragraph beginning ``The known band-independence exception, the Bashyal''}
% From app_statmethods.tex:446 @d55f7aad6; qualifier detail or procedure (DECISION-20261120a, second pass)
% [lead-in, Option B] In the paragraph as the note printed it, this text followed: ``The quoted 2D covariance therefore remains the block sum, with the cross term as a diagnostic.''\par
A fitted $(\Phi_{\text{norm}},\Phi_{\text{shape}},\text{MuonE}_{\text{scale}})$ posterior from MAT-MINERvA could supply a paper-independent $\rho$. The collaboration uses the same truncated-spectral pseudo-inverse as $\chi^{2}_{\ours}$ (private communication), so the convention is reported to be collaboration practice rather than one introduced by this analysis. No published reference for it has been identified.

\paragraph{MAT conformance of the band sum.}
% [records reorg] former heading: \subsection{From App. A, appstatmethods.tex: the paragraph beginning ``The combined chi2/mathrmndf=chiCombined is the descriptive''}
% From app_statmethods.tex:451 @d55f7aad6; qualifier detail or procedure (DECISION-20261120a, second pass)
The per-band sum is MAT-conformant, with the Flux band weighted as MAT weights it (\S\ref{sec:minerva}) and the $1.4\%$ normalization as a rank-1 band.

\paragraph{Toy pulls and training-seed block.}
% [records reorg] former heading: \subsection{From App. A, appstatmethods.tex: the paragraph beginning ``Across 205 reported bins, combined OmniFold''}
% From app_statmethods.tex:491 @d55f7aad6; qualifier detail or procedure (DECISION-20261120a, second pass)
% [lead-in, Option B] In the paragraph as the note printed it, this text followed: ``Across $205$ reported bins, combined OmniFold $+$ paper uncertainties give $\chi^{2}/\mathrm{ndf} = \chiCombined$ and pull mean/standard deviation $0.034/0.389$. This distance is descriptive, without calibrated goodness-of-fit (\S\ref{sec:chisq}).''\par
Same-ensemble toy pulls are approximately Gaussian; the $0.048\%$ training-seed block is more than an order of magnitude below the systematic block.

\paragraph{Edge-bin sensitivity.}
% [records reorg] former heading: \subsubsection{From App. A, the edge-bin sensitivity (row n81)}
% From app_statmethods.tex:358 @e89694209; register row, check stage 2026-11-21
The comparison recomputes $\chi^{2}/\mathrm{ndf}$ on shrinking subsets using only the paper covariance, with no OmniFold addition, for the exact-split cross-check (not recomputed for the LightGBM central value); the table is in the records document (Ref.~\cite{NoteRecords2026}).
The residual is concentrated in the low-$\pPar$, near-edge cells; the high-$\pPar$ tails carry sub-$2\%$ shape disagreement. This is a robustness check, not a separate method.

\paragraph{Fixed-truth coverage rescoring.}
% [records reorg] former heading: \subsubsection{From App. A, the fixed-truth coverage rescoring}
% From sec_validation.tex:102 @e89694209; check-stage claims read, second round
The rescoring is post hoc and partly circular, because the band's MC part reuses the toys' resampling mechanism.

\paragraph{Injected $W$ bump in five dimensions.}
% [records reorg] former heading: \paragraph{From the note's Table~\ref{tab:test-index}.}
% From part5_validation.tex (draft de1543818)
[Injected $W$ bump (5D), Outcome cell, after ``pass, by a narrow margin'':] Neighboring $W$ bins \SI{3.1}{\percent} and \SI{2.7}{\percent} high; per-5D-bin spread \SI{3.7}{\percent}

\subsection{Multidimensional statistical methods}\label{recsec:methods-nd}

\paragraph{Scope of the multidimensional appendix.}
% [records reorg] former heading: \subsection{From App. A, appstatmethodsnd.tex: the paragraph beginning ``This appendix gives the conventions for''}
% From app_statmethods_nd.tex:5 @d55f7aad6; qualifier detail or procedure (DECISION-20261120a, second pass)
This appendix gives the conventions for covariances and comparisons beyond the 2D result: rank, ndf and finite-ensemble declarations for N-D covariances, the scalar five-dimensional diagnostic reductions and interval definitions, the extended-fiducial band and pull fractions, and inference from unfolded events and selected projections.

\paragraph{Rank deficiency of N-D covariances.}
% [records reorg] former heading: \subsection{From App. A, appstatmethodsnd.tex: the paragraph beginning ``The 2D justification in Srefsec:chisq does''}
% From app_statmethods_nd.tex:10 @d55f7aad6; qualifier detail or procedure (DECISION-20261120a, second pass)
% [lead-in, Option B] In the paragraph as the note printed it, this text followed: ``An earlier 5D covariance from the standard production chain has $10\,694$ reported bins, effective positive rank $\mathbf{263}$ and infinite numerical condition number. It is distinct from $C_Z$, the five-axis covariance of the \resFiveAxis{} (Table~\ref{tab:config-key}), and was built before $C_Z$ was chosen for use. Its $4825$-bin 5D$\to$4D projection retains rank $263$ exactly.''\par
The covariance is a sum of about 45 rank-one outer products, one for each two-endpoint MAT systematic band, and of the statistical and machine-learning (ML) sample blocks. This construction cannot span $10\,694$ directions and is rank deficient. Quoting $\chi^{2}/n_{\text{reported}}$ would understate tension by roughly $40\times$ in 5D or $18\times$ in 4D.
% [records reorg] former heading: \subsection{From App. A, appstatmethodsnd.tex: the paragraph beginning ``For an N-D covariance, mathrmndf is''}
% From app_statmethods_nd.tex:12 @d55f7aad6; qualifier detail or procedure (DECISION-20261120a, second pass)
% [lead-in, Option B] In the paragraph as the note printed it, this text followed: ``For an N-D covariance, $\mathrm{ndf}$ is the retained rank.''\par
The reported-bin count $n_{\text{reported}}$ does not set it.
% [records reorg] former heading: \subsection{From App. A, appstatmethodsnd.tex: the paragraph beginning ``For an N-D covariance, mathrmndf is''}
% From app_statmethods_nd.tex:12 @d55f7aad6; qualifier detail or procedure (DECISION-20261120a, second pass)
In the paragraph as the note printed it, this text followed: ``Where an N-D $\chi^{2}$ is quoted, it must state five things. They are (i) the inverse actually used, including the pseudo-inverse cutoff \texttt{rcond} (singular values below this fraction of the largest are discarded) or the truncation rank, (ii) the retained rank, as $\mathrm{ndf}$, and (iii) the rank-truncation scan of that covariance.''\par
That scan is expected to lack the smoothness that supports the 2D choice in \S\ref{sec:chisq}.

\paragraph{Recording of ensemble sizes.}
% [records reorg] former heading: \subsection{From App. A, appstatmethodsnd.tex: the paragraph beginning ``Item (v) records the ensemble size''}
% From app_statmethods_nd.tex:16 @d55f7aad6; qualifier detail or procedure (DECISION-20261120a, second pass)
Item (v) records the ensemble size of each block. Without it, an uncorrected finite-ensemble bias could enter unnoticed.
% [records reorg] former heading: \subsection{From App. A, appstatmethodsnd.tex: the paragraph beginning ``Item (v) records the ensemble size''}
% From app_statmethods_nd.tex:16 @d55f7aad6; qualifier detail or procedure (DECISION-20261120a, second pass)
% [lead-in, Option B] In the paragraph as the note printed it, this text followed: ``For this analysis's production convention, universe-mean centered with biased $1/N$ normalization, it is $(N-p-2)/N$.''\par
For fixed residuals an uncorrected bias inflates $\chi^{2}$ and overstates tension. Confidence regions from the same over-tight covariance become over-optimistic.

\paragraph{Finite-ensemble debiasing.}
% [records reorg] former heading: \subsection{From App. A, appstatmethodsnd.tex: the paragraph beginning ``No generic correction is prescribed. A''}
% From app_statmethods_nd.tex:19 @d55f7aad6; qualifier detail or procedure (DECISION-20261120a, second pass)
% [lead-in, Option B] In the paragraph as the note printed it, this text followed: ``No generic correction is prescribed.''\par
A sum of deterministic rank-one band terms and statistical and ML sample blocks, such as the earlier 5D covariance from the standard production chain, has no single debiasing factor.
% [records reorg] former heading: \subsection{From App. A, appstatmethodsnd.tex: the paragraph beginning ``No generic correction is prescribed. A''}
% From app_statmethods_nd.tex:19 @d55f7aad6; qualifier detail or procedure (DECISION-20261120a, second pass)
% [lead-in, Option B] In the paragraph as the note printed it, this text followed: ``The standard factor also assumes independent Gaussian realizations and a data-independent truncation dimension, neither established for a data-dependent cut.''\par
The appropriate treatment depends on the covariance and inverse actually used, and is to be determined from the quantities recorded under item (v).

\paragraph{Normal-model approximations.}
% [records reorg] former heading: \subsection{From App. A, appstatmethodsnd.tex: the paragraph beginning ``For experiment i and cell f,''}
% From app_statmethods_nd.tex:33 @d55f7aad6; qualifier detail or procedure (DECISION-20261120a, second pass)
In the paragraph as the note printed it, this text followed: ``SD intervals use $s_f\sqrt{(n-1)/\chi^2_{0.975,n-1}}$ and $s_f\sqrt{(n-1)/\chi^2_{0.025,n-1}}$.''\par
These normal-model approximations condition on finite simulation and make no simultaneous claim across cells. A deterministic noise-free trace from one simulation receives no error band. Its cell-to-cell variation is not sampling uncertainty.

\paragraph{Interval assessment examples.}
% [records reorg] former heading: \subsection{From App. A, appstatmethodsnd.tex: the paragraph beginning ``The development assessment includes an interval''}
% From app_statmethods_nd.tex:35 @d55f7aad6; qualifier detail or procedure (DECISION-20261120a, second pass)
In the paragraph as the note printed it, this text followed: ``The table gives integer hits, both empirical fractions, exact pointwise 95\% binomial intervals and fixed half-width relative to mean truth.''\par
It conditions on estimated bootstrap widths without propagating finite-replica uncertainty.
% [records reorg] former heading: \subsection{From App. A, appstatmethodsnd.tex: the paragraph beginning ``The development assessment includes an interval''}
% From app_statmethods_nd.tex:35 @d55f7aad6; qualifier detail or procedure (DECISION-20261120a, second pass)
% [lead-in, Option B] In the paragraph as the note printed it, this text followed: ``Coverage fractions are not pooled across cells to certify them.''\par
The examples shown are a populated low-energy cell (EW7) and the corner cell EW41, already known to fail. They are not freshly selected validation targets.

\paragraph{Diagnostic reductions and checkpoints.}
% [records reorg] former heading: \subsection{From App. A, appstatmethodsnd.tex: the paragraph beginning ``For prior vertex v, data sensitivity''}
% From app_statmethods_nd.tex:47 @d55f7aad6; qualifier detail or procedure (DECISION-20261120a, second pass)
In the paragraph as the note printed it, this text followed: ``For prior vertex $v$, data sensitivity is $d_{vf}=\hat f_{{\rm data},v}/\hat f_{{\rm data},0}-1$, with envelope $h_f=\max_v|d_{vf}|$.''\par
Its denominator differs from the simulated residual's.
% [records reorg] former heading: \subsection{From App. A, appstatmethodsnd.tex: the paragraph beginning ``For prior vertex v, data sensitivity''}
% From app_statmethods_nd.tex:47 @d55f7aad6; qualifier detail or procedure (DECISION-20261120a, second pass)
% [lead-in, Option B] In the paragraph as the note printed it, this text followed: ``Neither $h_f$ nor an outer product of closure residuals is given a probability model.''\par
The reductions keep the per-cell values, the identity of each input and the exclusions, and distinguish complete outputs from checkpoints that carry no metadata. From a checkpoint, only the iterations that reproduce the originally recorded values are used. Later entries are excluded. The analyst cross-checked the reductions numerically.

\paragraph{Inference from unfolded events.}
% [records reorg] former heading: \subsection{From App. A, appstatmethodsnd.tex: the paragraph beginning ``Unfolding correlates weighted events. Treating them''}
% From app_statmethods_nd.tex:58 @d55f7aad6; qualifier detail or procedure (DECISION-20261120a, second pass)
% [lead-in, Option B] In the paragraph as the note printed it, this text followed: ``Treating them as independent in a downstream likelihood can underestimate parameter uncertainties~\cite{Desai:2025correlated}.''\par
Nominal weights do not specify uncertainty. Coherent unfolding replicas could propagate it through the fit. Resampling unfolded rows independently is not equivalent to resampling inputs and repeating the learned estimator. Classifier and permutation goodness-of-fit tests also require null calibration for their fitting and data reuse. These are requirements on such inference, with no new validation result here.
% [records reorg] former heading: \subsection{From App. A, appstatmethodsnd.tex: the paragraph beginning ``The uncertainty target matters: separate bin''}
% From app_statmethods_nd.tex:60 @d55f7aad6; qualifier detail or procedure (DECISION-20261120a, second pass)
% [lead-in, Option B] In the paragraph as the note printed it, this text followed: ``A fixed linear projection $y=P x$ propagates covariance as $C_y=P C_x P^{\top}$, with integrated bin widths in $P$ for densities.''\par
The identity neither calibrates that covariance nor accounts for regions chosen after inspecting data.

\subsection{Calibrated joint generator tests}\label{recsec:methods-joint}

% Copied whole and verbatim from sec_joint.tex @ b57c897b7 in the compression under DECISION-20261120a: each paragraph as the
% note printed it, with its comment lines; the note keeps the sentences a reader needs and points here. move-manifest-je JE-21.
% [records reorg] former heading: \subsection{Calibrated joint generator tests: the full paragraphs the note shortens}

\paragraph{Opening.}
% [records reorg] former heading: \subsubsection{Opening}
Two properties must be distinguished.  The coverage of unfolded intervals is the rate at which they contain the true cross section.  The calibration of a test is the rate at which it rejects a true hypothesis.  Calibration can be obtained without validated intervals by simulating the whole analysis chain at the hypothesis, so that the estimator's bias enters the simulated null distribution.  A calibration obtained in this way is valid to the extent that the simulated chain reproduces the analysis of the data.  Unfolded results can be compared with any prediction at truth level without simulating each generator through the detector. For hypothesis tests the advantage is smaller.  A test at reconstruction level would avoid the regularization bias, but it would still need the same response model.  Each external prediction would have to be folded through the analysis simulation's response, which is the hybrid construction of \S\ref{sec:joint5d-design}, with the same composition below the fine grid.  The tests of the note's joint-test section (\S\ref{sec:joint5d}) instead calibrate the unfolded statistic in simulation, so that the bias enters the null distribution rather than the decision.  We have not compared their power with that of a reconstruction-level test.

\paragraph{Hypotheses, grid and statistics.}
% [records reorg] former heading: \subsubsection{Hypotheses, grid and statistics}
For each of five predictions $G$ (the MINERvA Tune~v1, GENIE~2.12.10 without and with MEC, NuWro~21.09 and GiBUU~2019; \S\ref{sec:eavailw-band}), the null $H_{0}(G)$ is that the true cross section equals $G$ on every cell of a fine five-dimensional grid. Below that grid, the null truth has the MINERvA Tune~v1 simulation's event composition, varied in each pseudo-experiment by the drawn interaction-model universes. Every fine-cell integral is held fixed in that variation. The background is the Tune~v1 simulation's, with its uncertainty bands. This \emph{hybrid null} is the hypothesis the calibration simulates. For the Tune~v1, which is the analysis simulation itself, no external prediction is combined with the Tune~v1 composition. Three requirements are distinct. The hybrid null is defined exactly on the finite grid, and no convergence is needed to define it. Whether it represents the physical generator prediction depends on condition~(i) (\S\ref{sec:joint5d-conditions}). Whether the calibration reproduces the analysis of the data is calibration fidelity (conditions~(ii)--(iv)). Each prediction is tested in total (rate and shape together) and in shape alone: \jtNtests{} tests. The fine grid is the analysis binning of \S\ref{sec:method-binning}, $14\times16\times7\times7\times6$ bins in $(\pT,\pPar,\Eavail,q_{3},W)$; the test cells below are its integrals.
% Methods-paper shortening pass (2026-10-21, DECISION-20261020-joseph-refined-estimator-shorten-delta.md item 2): the null restated as in the article (its M8 correction,
% M2-AUDIT.md Q1 text fix, :46, as corrected by M2-ARCHIVE-CHECKS.md:6-8; code nd-unfolding/s5p_nullexp.py:299, :158-174,
% :215, :213; s5n_pseudo.py:135-142). Replaces "with the Tune v1's shape and G's finite-sample fluctuations below that
% grid" and "for the Tune v1 ... it is H0 exactly".
% plain-language pass 2 (DECISION-20261018-joseph-plain-language-pass2): pass-1 paper F7, note home first.
% Source: nd-unfolding/s5p_stage1_inspect.py:38-44 (FINE, the edges of sec_method.tex:99-101, beside J_EDGES,
% the coarse test grid below); RC4 fine-to-J map U (109, 65856), 65856 = 14*16*7*7*6,
% docs/publication/audit-20261008/notes/mapping-notes-C.md:25.
The tests are evaluated on a coarse $3^{5}$ grid in $(\pT,\pPar,\Eavail,q_{3},W)$ with edges $\pT\in[0,0.55,0.85,4.5]\,\si{GeV/c}$, $\pPar\in[1.5,3.5,6,60]\,\si{GeV/c}$, $\Eavail\in[0,0.4,1.5,100]\,\si{GeV}$, $q_{3}\in[0,1.2,2,100]\,\si{GeV}$ and $W\in[0,1.4,2.2,100]\,\si{GeV}$. Each inner edge is the fine-grid edge nearest a tertile of the weighted simulated truth distribution of the signal sample in that coordinate, so the cells keep physical bin widths. That distribution is taken from the development simulation. The tests are restricted to the \jtNcells{} supported cells (J). For GiBUU, whose sample lacks high-energy neutrinos, the domain is the \jtNcellsGiBUU{} cells with $\pPar<\SI{6}{GeV/c}$. The grid specifies a finite hybrid hypothesis and a reporting statistic. It does not turn OmniFold's event-level training into binned unfolding.
% Edge rule: nd-unfolding/s5c_partition.py:14-15 (docstring), :45-52 (coarse_edges: the interior fine edge nearest each
% weighted quantile k/parts, parts = 3 for H3), :80 (MCgen truth coordinates weighted by w_truth); development MC only
% (:2-7). No support threshold is applied in that file (it reports per-cell support only).
Two statistics are formed from the unfolded cell values $F$ and the prediction $\mu_{G}$. The total (rate-and-shape) statistic is $T_{\mathrm{tot}}=r^{\top}\Sigma^{-1}r$ with $r=F-\mu_{G}$. The shape statistic $T_{\mathrm{sh}}$ has the same form on the normalized cell fractions, with the normalization Jacobian applied to $\Sigma$ and one cell dropped. The metric $\Sigma=V+\mathrm{diag}\,\mathrm{Var}(\mu_{G})$ adds the prediction's own finite-sample variance to $V$. Here $V$ is a Ledoit--Wolf-shrunk covariance~\cite{LedoitWolf:2004} of \jtVn{} pilot pseudo-experiments at the Tune~v1 null. The metric defines how hypotheses are compared, and so it affects power. The null calibration below, not the metric, sets each test's size.
The prediction variance is diagonal because each simulated event of a prediction falls in one cell. The pilot's seeds are disjoint from those of the calibration, but it reuses the same simulated events and omitted the MINOS-efficiency band. The diagonal of $V$ matches the calibration's null variances (median ratio 1.00). The shrinkage, of intensity 0.015, floors about 80 of the 109 nearly degenerate correlation directions. At the Tune~v1 null the statistic therefore has a whitened trace of about 47 rather than 109. Since $V$ was fixed before the calibration, this affects power only (the measured powers are in the subsection ``Generator-test detail moved from the note'', above). The power at other shrinkage intensities has not been measured.
% Search: other shrinkage intensit; different lambda; shrinkage scan (M2-ARCHIVE-CHECKS.md:98-125 scans the trace only) @ 1a015eef

\paragraph{Calibration within the implemented simulation model.}
% [records reorg] former heading: \subsubsection{Calibration within the implemented simulation model}
The null distribution of each statistic is obtained from pseudo-experiments that run the full analysis chain at the null truth, using the joint-test estimator (\S\ref{sec:val-departure}). That estimator is deterministic LightGBM OmniFold with five iterations and the negative-weight background treatment, whose refinement classifier has 400 trees and 31 leaves. Each pseudo-experiment
\begin{itemize}[itemsep=1pt, topsep=2pt]
  \item draws Poisson pseudo-data from one half of the simulation together with an independent background source, refits the background-refinement classifier, and unfolds with the estimator trained on the other half;
  \item draws one of \jtFluxUniverses{} flux universes and every one of the \jtModelBands{} interaction-model bands at one of its universes;
  \item draws the weight-based detector bands (MINOS muon-matching efficiency and the GEANT4 neutron, pion and proton interaction reweights) and the five muon and beam-angle lateral bands through a linear surrogate;
  \item draws a \SI{1.4}{\percent} normalization, the data's rounding noise (taken from 20 real-data re-unfolds with rounding-level input changes), and the prediction's finite-sample residual.
\end{itemize}
The data are unfolded once, with the full simulation. The Monte Carlo $p$ value of a test is $(k+1)/(B+1)$, where $k$ of $B$ null pseudo-experiments reach the observed statistic. This $p$ value is never zero. The $p$ value on which each decision rests, the \emph{claim} $p$ value, is the largest over the specified \emph{claim variants}, which shift the null. The variants move it by an upper bound on the simulation-size difference between the pseudo-experiments and the data unfold (condition~(iv), \S\ref{sec:joint5d-conditions}), and by $\pm\jtMoneKappa\,\delta_{\mathrm{fine}}$ (condition~(i)). The external nulls are strongly non-central, with non-centralities $\lambda$ of \jtLambdaMin{}--\jtLambdaMax{} in these units. They are non-central because the calibration reproduces the estimator's pull toward its prior (\S\ref{sec:val-departure}). Each observed statistic is compared with that displaced distribution (the earlier note's figure of the joint tests' null distributions).

\paragraph{Confidence intervals for functionals.}
% [records reorg] former heading: \subsection{From App. A, appstatmethodsnd.tex: the paragraph beginning ``Recent work constructs confidence intervals for''}
% From app_statmethods_nd.tex:62 @d55f7aad6; qualifier detail or procedure (DECISION-20261120a, second pass)
Recent work constructs confidence intervals for selected functionals in constrained inverse problems~\cite{Stanley:2025functionals} and joint confidence regions for multiple functionals~\cite{Batlle:2025simultaneous}, which could provide independent benchmarks for specified integrals or contrasts. Their guarantees depend on the forward model, noise law, and parameter constraints and do not transfer automatically to a learned response with finite simulation and detector nuisances. A rank-deficient forward operator is also distinct from a rank-deficient sample covariance. Neither truncating the latter nor validating a few projections certifies the full distribution. Such a benchmark would supplement the scalar five-dimensional uncertainty construction that a measurement would require. It would not replace that construction.

\paragraph{Empty cells in calibration pseudo-experiments.}
% [records reorg] former heading: \subsection{From App. A, appstatmethodsnd.tex: the paragraph beginning ``Each calibration pseudo-experiment splits the simulation''}
% From app_statmethods_nd.tex:69 @d55f7aad6; qualifier detail or procedure (DECISION-20261120a, second pass)
% [lead-in, Option B] In the paragraph as the note printed it, this text followed: ``\paragraph{Generative model of one pseudo-experiment.} Each calibration pseudo-experiment splits the simulation at random into two halves (Table~\ref{tab:joint5d-genmodel}) and refits the background refinement. Its truth equals $G$ on every fine cell that the source half populates.''\par
Cells it leaves empty lower a coarse cell by at most \SI{0.2}{\percent}, or about \SI{4}{\percent} of the null standard deviation (the per-experiment spread is below \SI{1.5}{\percent} of it).

\paragraph{Monte Carlo independence of predictions.}
% [records reorg] former heading: \subsection{From App. A, appstatmethodsnd.tex: the paragraph beginning ``The data's rounding noise, and the''}
% From app_statmethods_nd.tex:71 @d55f7aad6; qualifier detail or procedure (DECISION-20261120a, second pass)
% [lead-in, Option B] In the paragraph as the note printed it, this text followed: ``The prediction's variance in that draw is its own per-cell sum of squared weights. That diagonal model is an assumption of the metric.''\par
Disjoint cells do not by themselves make a prediction's Monte Carlo errors independent. The background source and its template are Poisson draws from disjoint halves of the background sample. The detector weight bands include the GEANT4 hadron-interaction reweights.

\paragraph{Lost pseudo-experiments and test size.}
% [records reorg] former heading: \subsection{From App. A, appstatmethodsnd.tex: the paragraph beginning ``emphProbability model. Fix one test, and''}
% From app_statmethods_nd.tex:132 @d55f7aad6; qualifier detail or procedure (DECISION-20261120a, second pass)
In the paragraph as the note printed it, this text followed: ``Assumption A2 is that, given the data, these are independent Bernoulli draws with a common probability $p$, the calibrated tail probability of the observed statistic.''\par
This requires the lost pseudo-experiments to be lost independently of their outcomes.
% [records reorg] former heading: \subsection{From App. A, appstatmethodsnd.tex: the paragraph beginning ``emphProbability model. Fix one test, and''}
% From app_statmethods_nd.tex:132 @d55f7aad6; qualifier detail or procedure (DECISION-20261120a, second pass)
% [lead-in, Option B] In the paragraph as the note printed it, this text followed: ``Assumption A1 is that, for a true null, $p$ is super-uniform over repetitions of the experiment, $P(p\le u)\le u$, with the nuisance parameters drawn from their priors.''\par
The size statement is therefore averaged over the nuisance priors.

\paragraph{Multiplicity, stopping rule and design.}
% [records reorg] former heading: \subsubsection{Multiplicity, stopping rule and design}
The ten tests are carried out with the Holm procedure~\cite{Holm:1979} at familywise level $\alpha=\jtAlpha$. A determinacy requirement also applies: a step rejects only when the 95\% Clopper--Pearson interval~\cite{Clopper:1934} of $k/B$ lies entirely below its Holm threshold. Calibration proceeded in batches of \jtBatch{} pseudo-experiments per null, to at most \jtBcap{}, under a sequential rule that stopped a null once both of its decisions were determined. A null stopped when a 99.5\% Clopper--Pearson interval of $k/B$ lay entirely on one side of each of its thresholds and met a fixed precision requirement. The decision rule uses the upper end of the two-sided 95\% interval, a one-sided 97.5\% upper bound.
Two assumptions are made: (A1) each true null's calibrated tail probability is uniformly distributed or conservative, and (A2) the pseudo-experiments' exceedances are independent given the data. Under A1 and A2, the probability that this upper end falls below the threshold, at any number of pseudo-experiments up to \jtBcap{}, is below the threshold. This holds whatever rule determines when sampling stops, including the budget extensions. The decision rule therefore controls the familywise error rate of the ten decisions at 5\% under these assumptions. That rate is the probability of falsely rejecting at least one true hypothesis among them. The error control applies only to the decisions. The point estimates $(k+1)/(B+1)$ carry no separate error-control guarantee. Under these assumptions, the bound $\mathrm{FWER}\le0.8902\,\alpha<0.9\,\alpha$ is derived in \S\ref{sec:joint5d-fwer}.
A2 requires the lost pseudo-experiments to have been lost independently of their outcomes, which has not been shown. When rerun, however, none of them reaches the observed statistic (\S\ref{sec:joint5d-lost}). The calibrated size is averaged over the nuisance priors: at a fixed true value of the nuisance parameters it can exceed $\alpha$.
% Methods-paper shortening pass (2026-10-21, DECISION-20261020-joseph-refined-estimator-shorten-delta.md item 2): moved from the article with its sources.
% Stopping: a 99.5% look interval plus a precision tier (nd-unfolding/s5p_inference.py:89-114, look_level=0.995), the
% decision the 95% interval (:127-155). Gate and bound: docs/orchestration/state/2d-closeout/stageF/methods/
% M2-FWER-RECORD.md:69-74 (text, provisional J4), :49-52 (bound), :56-66 (assumptions A1, A2 and the rule); X5 G9: the
% one- vs two-sided wording (m2-calc/certify_crossing.py:1-2, U95 = beta.ppf(0.975, k+1, B-k)); nuisance-averaged size
% M2-AUDIT.md:143 (Q3).
% Search: lost independently of their outcomes; outcome-independent; independent of their outcomes @ ddf16f10
The joint-test estimator (\S\ref{sec:joint5d-calibration}) was chosen before the design was fixed. The hypotheses, statistics, variants, multiplicity rule, stopping rule and every input were fixed before any production pseudo-experiment or observed statistic, with two data-informed elements. First, the design's reviewers had seen descriptive per-cell residuals of the data before it was fixed. Second, a bounded extension of the production budget was decided during production, while the calibration's monitoring output showed observed claim $p$ values, so that the sequential rule, rather than the budget, could set the stop. No statistic, variant, threshold or stopping rule was changed, and the familywise bound above holds whatever rule determined when sampling stopped. An independent recomputation from the stored products, with separately written code, reproduces every $p$ value and decision of the production evaluation. It verifies the calculation, not the adequacy of the calibration. The calibration is valid for the simulated null to the extent that the pseudo-experiment chain reproduces the analysis applied to the data. No rejection rate has been measured for null pseudo-data processed like the data, which would require an independent simulation sample.
% Methods-paper shortening pass (2026-10-21, DECISION-20261020-joseph-refined-estimator-shorten-delta.md item 2): the article's M8 replacement of "the Monte Carlo
% p values remain valid at the number of pseudo-experiments reached" (M2-FWER-RECORD.md:69-74), applied here.
% Register round n1 (2026-10-25): the calibration's validity statement and the unmeasured rejection rate moved here
% from the calibration-validity paragraph of sec:joint5d-conditions (M8; M2-AUDIT.md Q3, Q4e, :187-193).
% Search: rejection rate; processed like the data; independent simulation sample @ ddf16f10
Three studies followed the decisions: the scan of the shift multiplier (\S\ref{sec:joint5d-kappa}), the rerun of the lost pseudo-experiments (\S\ref{sec:joint5d-lost}), and the matched coarse-projection comparison (\S\ref{sec:joint5d-projections}). They are for reporting only, outside the decision rule. Each was specified before it ran. The projection comparison was specified after the five-dimensional outcomes were known.

\paragraph{Holm reduction without independence.}
% [records reorg] former heading: \subsection{From App. A, appstatmethodsnd.tex: the paragraph beginning ``emphHolm reduction without independence. Let m0''}
% From app_statmethods_nd.tex:140 @d55f7aad6; qualifier detail or procedure (DECISION-20261120a, second pass)
% [lead-in, Option B] In the paragraph as the note printed it, this text followed: ``A familywise error therefore requires some true null to cross below $\alpha/m_{0}$, and by the union bound $\mathrm{FWER}\le m_{0}\,R(\alpha/m_{0})\le\max_{j\le m}j\,R(\alpha/j)$.''\par
The argument uses neither independence between tests nor the order of the Holm steps, so the shared ensemble of a null's total and shape tests, their unequal stopping counts, and the ranking by $(k+1)/(B+1)$ with its ties do not affect it.

\paragraph{Familywise error: the crossing event and the worst case.}
% [records reorg] former heading: \subsection{Familywise error: the crossing event and the worst case (from App. A, sec:joint5d-fwer)}
% From app_statmethods_nd.tex:134 @e57f0a0a8; derivation steps
\emph{Crossing event.}  The rule rejects at threshold $c$ only if $U(k_{B},B)<c$ at the number $B\le\jtBcap$ at which sampling stopped.  Whatever rule chose that stopping point, including stops decided with results visible, irregular looks, lost draws, and the budget extensions, the rejection lies in the event $E_{c}=\bigcup_{B\le\jtBcap}\{U(k_{B},B)<c\}$ that the bound crosses below $c$ at some $B$.  Hence $P(\text{reject at }c)\le R(c)\equiv E_{p}[g_{c}(p)]$, with $g_{c}(p)=P(E_{c}\mid p)$.
\emph{Uniform $p$ is the worst case.}  Coupling the Bernoulli draws in $p$ shows that $g_{c}(p)$ does not increase with $p$.  For any super-uniform $p$ its expectation is therefore at most its value for uniform $p$, $R(c)\le\int_{0}^{1}g_{c}(u)\,du$.

\paragraph{Familywise error: certified evaluation of $R(c)$.}
% [records reorg] former heading: \subsection{Familywise error: certified evaluation of R(c) (from App. A, sec:joint5d-fwer)}
% From app_statmethods_nd.tex:145 @e57f0a0a8; derivation steps
\emph{Certified constant.}  For each $c=\alpha/j$, $j=1,\dots,10$, let $k^{*}(B)$ be the largest $k$ with $U(k,B)<c$.  For fixed $p$, a dynamic program propagates the distribution of $k_{B}$ from $B=1$ to \jtBcap{}, moving the probability of every state with $k_{B}\le k^{*}(B)$ into an absorbing ``crossed'' state at each $B$, and pooling the counts above $\max_{B}k^{*}(B)$, which can never cross because counts never decrease. The absorbed probability is $g_{c}(p)$ exactly.
This is evaluated on a grid of 6000 points in $[0,6c]$.  Because $g_{c}$ is nonincreasing, a left Riemann sum over-estimates its integral on each grid interval, and with the tail term $(1-6c)\,g_{c}(6c)$ it is a rigorous upper bound on $\int_{0}^{1}g_{c}$, so the integration error is included in the bound.
% [records reorg] former heading: \subsection{From App. A, appstatmethodsnd.tex: the paragraph beginning ``The largest ratio R(c)/c, at j=1,''}
% From app_statmethods_nd.tex:142 @d55f7aad6; qualifier detail or procedure (DECISION-20261120a, second pass)
In the paragraph as the note printed it, this text followed: ``The largest ratio $R(c)/c$, at $j=1$, lies in $[0.8892,0.8902]$.''\par
It decreases to $[0.5138,0.5148]$ at $j=10$.

\paragraph{Decisions on the ten hypotheses.}
% [records reorg] former heading: \subsubsection{Decisions on the ten hypotheses}
Each decision holds under six conditions of the specified model (\S\ref{sec:joint5d-conditions}). The earlier note's table of the joint generator tests lists the ten decisions. Every test is rejected at familywise $\alpha=\jtAlpha$; nine have $k=0$, and NuWro's shape test has $k=1$. Each rejection also survives the robustness variants at $\pm\jtMoneKappaRobust\,\delta_{\mathrm{fine}}$, the upper end of the extrapolation in condition~(i). The earlier note's figure of the joint tests' null distributions shows each null's calibrated distribution with the observed statistic. The Tune~v1 null is centered at small values, whereas the external nulls are displaced far upward by the estimator's pull toward its prior.
The multiplier $\kappa$ scales the last measured refinement change $\delta_{\mathrm{fine}}$: $\kappa=1$ means one such change, not one standard deviation or a confidence level. A scan of $\kappa$ along the direction $\delta_{\mathrm{fine}}$ of the claim variants~\cite{JointTestKappa2026} (for reporting only) found, in the calibration ensembles used for the decisions, all ten decisions to remain determinate rejections at every sampled multiplier up to 5.25. The scan probes one prescribed direction. It establishes neither convergence nor a bound on the actual below-grid residual (condition~(i), \S\ref{sec:joint5d-conditions}). Where determinate rejection is first lost, and by which tests, is given in the records document.
Because every test rejects, power serves here only as context: the test specification uses power only to interpret a failure to reject. The powers measured at the Tune~v1 and GENIE~2.12.10 nulls are given in the records document.

\paragraph{Scan of the multiplier $\kappa$.}
% Moved verbatim from the note's generator tests (sec_joint.tex) in the compression (G2b; DECISION-20261119c item 3 and the director's rule (5),
% 2026-11-19: the later reporting-only studies' detail; the note keeps each study's conclusion and qualifiers; move-manifest-je JE-12).
% [records reorg] former heading: \subsection{Generator-test detail moved from the note}
% [records reorg, repeated sentence not printed] The multiplier $\kappa$ scales the last measured refinement change $\delta_{\mathrm{fine}}$: $\kappa=1$ means one such change, not one standard deviation or a confidence level.
% [records reorg, repeated sentence not printed] A scan of $\kappa$ along the direction $\delta_{\mathrm{fine}}$ of the claim variants~\cite{JointTestKappa2026} (for reporting only) found, in the calibration ensembles used for the decisions, all ten decisions to remain determinate rejections at every sampled multiplier up to 5.25.
% [records reorg, repeated sentence not printed] The scan probes one prescribed direction.
% [records reorg, repeated sentence not printed] It establishes neither convergence nor a bound on the actual below-grid residual (condition~(i), \S\ref{sec:joint5d-conditions}).
 Determinate rejection is first lost between the sampled multipliers 5.2547 and 5.2563, and between 5.1688 and 5.1703 with the recovered draws included. It is lost by the GENIE~2.12.10 + MEC shape test, whose interval then contains its Holm threshold, and, through the Holm step-down, by the GENIE~2.12.10 shape test. Neither is then a determinate failure to reject.
% Search: below the fine grid; sub-fine-grid; fine-grid convergence; delta_fine; kappa breakdown @ ddf16f10

\paragraph{Power of the joint tests.}
% [records reorg, repeated sentence not printed] Because every test rejects, power serves here only as context: the test specification uses power only to interpret a failure to reject.
 At the Tune~v1 null, two of the three specified alternatives are detected in shape with power 1.0 at the 0.005 level, and the third, a MEC-like $(\Eavail,W)$ shape, with \jtPowPthree. At the GENIE~2.12.10 null the shape powers at 0.005 are \jtPowPonegShape, 0.02 and 0.00. The binomial 95\% intervals are $[0.98,1]$ for the first two Tune~v1 alternatives, $[0.17,0.30]$ for the third, and $[0.67,0.80]$, $[0.01,0.05]$ and $[0,0.02]$ at the GENIE~2.12.10 null. These values use the power sets as evaluated for the decisions. Each of those sets had lost some of its 200 alternative pseudo-experiments to time limits (172--199 retained). Evaluated against the null ensembles used for the decisions, the complete power sets, recovered as in \S\ref{sec:joint5d-lost}, agree with them within 0.01. When the GENIE~2.12.10 null is also recomputed on all draws, the shape power at 0.005 against the first alternative at that null is 0.72.
% X7 (sol minor), moved from the article: docs/orchestration/state/s5p/stage7/joint/joint-evaluate.json
% power.<P>.shape["0.005"].claim_rule_determined.interval (n 193, 195, 195; 199, 193, 172).

\paragraph{Six conditions attached to the ten rejections.}
% [records reorg] former heading: \subsubsection{Six conditions attached to the ten rejections}
Each decision holds under six conditions of the specified model. They are restated here in the terms of this note, numbered (i)--(vi) as in the test specification, and fall into three kinds.
% Deposit pass / r22 Step 2 (B10): the six conditions of RECORD-20261005-s5p-joint-5d-inference-result.md sec. 3
% (amendment 7 claims.conditions_stated_with_every_claim), formerly quoted verbatim here, are restated in prose; the
% verbatim text is in that record (app_provenance.tex, joint-5d). Content substitutions kept from the quoted version:
% (i) the null truth's below-grid description as implemented (M8 correction, M2-AUDIT.md Q1; sec:joint5d-design);
% (ii) "(app_statmethods.tex)" -> App. A; (iii) the clipping's effect measured on the pilot
% (EXCEPTION-20260927-s5p-pilot-negative-rate-resubmission.md, "Quantified effect": 0.35% of the band's shift in
% whitened norm at fine reco resolution); (iv) "~0.05 null SD" superseded by the frozen-V prefreeze measurement, max
% 0.068 (CLAIM-TABLE row J12); (v) "KNOWN_ISSUES 82" -> the GiBUU sample's energy limit, sec:eavailw-band.

\paragraph{Numerical convergence.}
% [records reorg] former heading: \paragraph{Numerical convergence.}
 Condition~(i) concerns the null truth below the fine grid. There it has the Tune~v1 composition of the hypothesis definition (\S\ref{sec:joint5d-design}), varied by the drawn interaction-model universes, with every fine-cell integral held fixed at that of $G$, including the finite-sample fluctuations of $G$'s simulation. The fine-grid convergence check has not shown convergence. In the $L_{2}$ norm of the per-cell relative changes, the last refinement, from a twofold-merged grid to the fine grid, is \jtMoneMin--\jtMoneMax{} of the whole change from the coarse test grid to the fine grid. In that norm the successive refinement changes are therefore not shown to shrink, and the residual below the fine grid is unmeasured and potentially material.
% Search: below the fine grid; fine-grid convergence; delta_fine; per halving @ ddf16f10
The design assumes that each twofold refinement of the grid multiplies the refinement change, measured in a proxy of the test metric, by a factor of \jtMoneWnormMin--\jtMoneWnormMax{}. This factor comes from the review of the design, and its derivation is not given. If the change kept shrinking geometrically at that rate, the residual would be about one to three times the last refinement change, $\delta_{\mathrm{fine}}$. Here $\delta_{\mathrm{fine}}$ is the change of the noise-free unfolded result when the null truth is refined from the twofold-merged grid to the fine grid. That estimate sets the claim variants at $\pm\jtMoneKappa\,\delta_{\mathrm{fine}}$, the central extrapolated estimate, and the robustness variants, for reporting only, at $\pm\jtMoneKappaRobust\,\delta_{\mathrm{fine}}$. The extrapolation assumes geometric convergence. The measured $L_{2}$ ratios also allow slower rates and larger residuals.
Measured with $V$ fixed before any observed statistic, the claim variants move the external nulls by \jtMoneClaimSdMin--\jtMoneClaimSdMax{} null standard deviations and the robustness variants by \jtMoneRobSdMin--\jtMoneRobSdMax{}. For the hybrid null as simulated, the residual plays no role. The variants carry the rejections over to $G$ itself only if replacing the Tune~v1's shape below the fine grid by $G$'s moves the noise-free unfolded result by no more than $\pm\jtMoneKappa\,\delta_{\mathrm{fine}}$ along $\delta_{\mathrm{fine}}$. Condition~(vi) concerns the predictions' own finite simulated samples, which enter the metric as $\mathrm{diag}\,\mathrm{Var}(\mu_{G})$ and the calibration as a drawn residual term.
% Search: below the fine grid; fine-grid convergence; delta_fine; per halving @ ddf16f10

\paragraph{Calibration validity.}
% [records reorg] former heading: \paragraph{Calibration validity.}
 The calibration's validity, and what the independent recomputation verifies, are stated in the records copy of the paragraph on the design, above. Two known departures are bounded. Condition~(iv) concerns the simulation size: the pseudo-experiments unfold with half of the simulation and the data with all of it. The resulting shift $D$, estimated from 16 pairs, is at noise level per cell, with a median $|D|/f$ of 0.02--0.075\%. Along the null's bias direction it is not always consistent with zero (proxy $a/\mathrm{se}$ of $+4.7$ for the Tune~v1 and $-3.4$ for NuWro), but it is at most 0.068 null standard deviations. That bound was measured with $V$ fixed before any observed statistic, and it supersedes the specification's estimate of about 0.05. The shift enters as the bias-aligned upper-bound variant. That shift was measured for the mean only, in a process without noise or nuisance draws.
A separate check reran the first 200 calibration pseudo-experiments of the NuWro, GENIE~2.12.10, and GENIE~2.12.10 + MEC nulls with a quarter of the simulation. Between half and quarter simulation, it resolved no change of the null distributions' location, width, or tails: the medians agree within 0.074 null standard deviations, and the widths lie within 0.975--1.010 of the half-simulation values. The only change it resolved is a mean paired shift of 0.006--0.034 null standard deviations. Under an assumed extrapolation from half to full simulation that is conservative by design, all three shape decisions remain rejected. With the margins of the other tests stated below, all ten remain rejected. The check concerns the simulation size only and is not a test of the data process.
Every pseudo-experiment resamples the same simulated events, so the null distributions are conditional on that simulation. The decisions closest to their boundary are the shape tests of NuWro and GENIE~2.12.10. With the Holm order held fixed, their observed statistics exceed the value at which the decision would change by 1.11 and 1.30 null standard deviations, and by 1.01 and 0.67 with the robustness variants at $\pm\jtMoneKappaRobust\,\delta_{\mathrm{fine}}$. The other eight tests lie 2.7 to 203 null standard deviations from that value, and 2.0 to 203 with the robustness variants. Crossing that value would make a test undetermined, its Clopper--Pearson interval straddling the threshold, rather than not rejected. The data unfold's rounding-level nondeterminism, measured over the 20 re-unfolds, leaves every observed $p$ value unchanged except NuWro's shape test (2/1752 to 3/1752).
% Methods-paper shortening pass (2026-10-21, DECISION-20261020-joseph-refined-estimator-shorten-delta.md item 2): "The Monte Carlo p values are exact for the
% simulated null" restated as the calibration's validity, as in the article (M8; M2-AUDIT.md Q3, Q4e, :187-193);
% register round n1 (2026-10-25): that statement now prints once, in sec:joint5d-frozen.
% Methods-paper shortening pass (2026-10-21, DECISION-20261020-joseph-refined-estimator-shorten-delta.md item 2): the article's simulation-size check (SF-1) and
% decision-boundary margins, moved with their sources. SF-1: docs/orchestration/state/2d-closeout/stageF/methods/
% SF1-REPORT.md:5, :23-35 (all 600 quarter products; |Delta| <= 0.074 null SD, s 0.975-1.010, P95/P99 ratios
% 0.987-1.008, every interval containing 0 or 1; mean paired shift 0.006-0.034 null SD, resolved; NuWro shape k = 2/1751
% after the most adverse adjustment, GENIE CV and MEC shape k = 0; all ten decisions rejected); method SF1-SPEC.md
% sec. 2-4; NOT CITABLE FOR the data process (SF1-REPORT.md:9). Margins:
% docs/orchestration/state/measurement-paper/phase1/a7margins/RECEIPT.md:65-78, :89-102, :55; a7margins.json f0960dc6...;
% descriptive, not a re-grade (:3).
% Search: test of the data process; quarter simulation; rejection rate; processed like the data @ ddf16f10
% 0.07: the frozen-V prefreeze measurement (max 0.068, GENIE+MEC), which supersedes amendment 7 condition 4's
% reviewer-proxy "~0.05" (docs/publication/audit-20261008/CLAIM-TABLE-20261008.tsv row J12;
% REVIEW-20261007-paper-completion.md E6). This section printed 0.05 until 2026-10-11.

\paragraph{Physical-model adequacy.}
% [records reorg] former heading: \paragraph{Physical-model adequacy.}
 Condition~(ii) concerns the detector model. It is the simulation's per event, varied only by the drawn detector bands (\S\ref{sec:joint5d-calibration}, each at $\pm1\sigma$) and by the linear surrogate for the five lateral bands. The surrogate has been checked only for the symmetry of its endpoints. The analysis's systematic set (App.~\ref{app:statmethods}) has no recoil-energy-scale band, so the completeness of the hadronic response is not established (the earlier note's section on the hadronic-response condition). Condition~(iii) concerns the interaction-model prior, the analysis's set of \jtModelBands{} interaction-model reweighting bands, with every band drawn in every pseudo-experiment. Negative weights in two of them, LowQ2\_1 and HighQ2\_1, are clipped per universe. Measured on the pilot, the effect of clipping them rather than keeping them is at most 0.35\% of a band's shift in whitened norm. Condition~(v) concerns the neutrino-energy range. The GiBUU sample contains no events with $E_{\nu}>\SI{20}{GeV}$ (\S\ref{sec:eavailw-band}), and its truth outside the test domain ($\pPar\ge\SI{6}{GeV/c}$, with truncated $E_{\nu}$) enters through the migrations into the domain cells. No prediction contains events with $E_{\nu}>\SI{100}{GeV}$; only the data's flux integral stops there, and the omission is negligible. The use of the Tune~v1 composition below the fine grid in (i) is also a modeling choice of this kind.
% M2-AUDIT.md Q4d, :183-184.
% Search: recoil-energy-scale band; calorimetric response band; linear surrogate; endpoint symmetry @ ddf16f10

\paragraph{Paired rounding diagnostic.}
% [records reorg] former heading: \subsubsection{From App. A, the paired rounding diagnostic}
% From app_statmethods_nd.tex:51 @e89694209; check-stage claims read, second round
This finite-sample diagnostic supports the construction without proving coverage or independence from every physical uncertainty.

\paragraph{Paired-model $\sqrt{2}$ factor.}
% [records reorg] former heading: \subsection{From App. A, appstatmethodsnd.tex: the paragraph beginning ``The rounding study keeps the fine-grid''}
% From app_statmethods_nd.tex:43 @d55f7aad6; qualifier detail or procedure (DECISION-20261120a, second pass)
% [lead-in, Option B] In the paragraph as the note printed it, this text followed: ``The estimated numerical share is $s_{\rm pair}^2/s_{\rm boot}^2$.''\par
The paired-model $\sqrt2$ factor introduces no physical noise source.

\paragraph{Lost pseudo-experiments and their recovery.}
% [records reorg] former heading: \subsubsection{Lost pseudo-experiments and their recovery}
Computing time limits interrupted or prevented \jtLostCal{} of the submitted calibration pseudo-experiments (\jtLostInterrupted{} interrupted, \jtLostNever{} never started) and \jtLostPow{} power pseudo-experiments. Every submitted draw has been observed, and no recovered draw reaches the observed statistic in any test, claim variant or robustness variant. The decisions in the earlier note's table of the joint generator tests remain those of the evaluation before recovery, at the $B$ at which its rule stopped (App.~\ref{app:devhistory-lost}).
% Search: earlier-stop; earlier stop; independent check of the earlier stops; standalone rerun @ 1a015eef
% r22 Step 2 (B10): the cross-check's count (856 of 856 quantities) is in app_provenance.tex, joint-5d.
% Deposit pass: the disposition label Joseph selected (DECISION-20261005-s5p-recompute-extension-and-lost-seed-
% recovery.md sec. 5) restated as its content, that section's "What it changes" and required disclosures 1-5,
% printed in the body summary and App. devhistory; old text in git history.
% Source: RECORD-20261006-s5p-lost-seed-recovery-resolution.md (recovery, margins 4.6/49.8 T units, earlier
% stops, power over the complete sets); RECORD-20261005 sec. 5 (pre-recovery sensitivity, |rho| <= 0.06);
% DECISION-20261005-s5p-recompute-extension-and-lost-seed-recovery.md sec. 5 (Joseph's disposition).

\paragraph{Evaluation of the recovered draws.}
% [records reorg] former heading: \subsubsection{Lost pseudo-experiments of the generator tests: evaluation of the recovered draws}
The recovered draws were evaluated in two ways: with the variant shifts of the primary evaluation held fixed (the \emph{fixed-shift} evaluation), and with the shifts recomputed on all draws (the \emph{recomputed-shift} evaluation). The recovery's checks are narrower in scope than the main recomputation. Only the fixed-shift evaluation was cross-checked independently. The recomputed-shift evaluation was reproduced only by a standalone rerun of the same code, a consistency check rather than an independent recomputation. The earlier stops rest on the calibration run's own re-application of the sequential rule and have no independent check.

\paragraph{Worst case and determinism check.}
% [records reorg] former heading: \subsubsection{Lost pseudo-experiments of the generator tests}
Had every interrupted draw exceeded the observed statistic, no test would have been rejected ($p$ \jtWorstIntLo--\jtWorstIntHi). Among completed draws, runtime is not associated with the statistic ($|\rho|\le\jtRuntimeRho$), which supports but does not certify the decisions.
A determinism check fixed before the rerun (16 of 16 completed seeds rerun bitwise identically) shows that each recovered draw is the lost draw, so every submitted draw has been observed. The closest margin is 4.6 units of $T$ below the observed value of 868.6 (GENIE~2.12.10 + MEC shape, $\pm\jtMoneKappaRobust\,\delta_{\mathrm{fine}}$ variant). The next is 49.8 (GENIE~2.12.10 shape, same variant). With complete batches the stopping rule, fixed before any production pseudo-experiment, would have stopped the Tune~v1, GENIE~2.12.10, GENIE~2.12.10 + MEC and GiBUU nulls one batch earlier ($B=1200$), with the same decisions.

\paragraph{Matched coarse projections.}
% [records reorg] former heading: \subsubsection{Matched coarse projections}
The same calibrated tests are restricted to matched $3\times3$ coarse projections of the same cells, onto $(\pT,\pPar)$ and onto $(\Eavail,W)$. The comparison uses the same calibration ensembles, with the recovered draws included, and the same variants and claim rule. Its $p$ values were reproduced independently~\cite{JointTestProjections2026}. The projected tests do not enter the Holm family and change no decision of the earlier note's table of the joint generator tests. The comparison and its criterion were specified after the five-dimensional outcomes were known but before the comparison was run.
Under the criterion, a test would count as evidence of additional joint discrimination only if both of its projected claim $p$ values were at least 0.05 in both the fixed-shift and the recomputed-shift evaluations. No test meets it, and the values are identical in the two evaluations. Every test has $p<0.05$, without a multiplicity correction, in at least one of the two projections. This study therefore provides no evidence of additional discrimination from the joint five-observable information. Table~\ref{tab:joint5d-projections} gives every projected claim $p$ value. In the muon-only $(\pT,\pPar)$ projection, GENIE~2.12.10 is not rejected ($p=\woneGcvPtTot$ in total, $\woneGcvPtShape$ in shape; a post hoc observation), whereas its $(\Eavail,W)$ projection is. Hadronic information therefore participates in that rejection. This does not establish added power of the joint test over all matched projections, or accept GENIE as true. The powers of the joint and projected tests have not been compared.
% PACKET-20261005-publication-decision.md sec. 6 (W1, "Criteria, fixed now"): "The marginal tests are report-only.
% They do not enter the Holm family and change no frozen decision." (docs/publication/, :252)
% Register round r12 (2026-10-21): moved from the article's Sec. VII E; source W1-RECORD-20261006.md (no power comparison).
% Search: power comparison; powers of the joint; no evidence of additional discrimination @ 1a015eef
% Deposit pass: W1-RECORD-20261006.md sec. 1 (identical under both D2 readings; independently reproduced, sec. 4, 40 of 40). Values copied as recorded.
\begin{table}[htbp]
  \caption[Matched-projection test probabilities]{Every test has $p<0.05$ in at least one matched projection. The joint tests are restricted to matched $3\times3$ coarse projections of the same cells (for GiBUU, 6 and 8 cells of its domain), with all submitted calibration draws, the same variants and the same claim rule. The GENIE~2.12.10 values in the $(\pT,\pPar)$ column are a post hoc observation.}
  \label{tab:joint5d-projections}
  \centering\small
  \begin{tabular}{@{}llll@{}}
    \toprule
    $G$ & test & $(\pT,\pPar)$: claim $p$ ($k/B$) & $(\Eavail,W)$: claim $p$ ($k/B$) \\
    \midrule
    MINERvA Tune~v1     & total & 0.00071 (0/1400) & 0.00071 (0/1400) \\
    MINERvA Tune~v1     & shape & 0.00071 (0/1400) & 0.00071 (0/1400) \\
    GENIE~2.12.10       & total & 0.483 (676/1400) & 0.00071 (0/1400) \\
    GENIE~2.12.10       & shape & 0.629 (880/1400) & 0.00071 (0/1400) \\
    GENIE~2.12.10 + MEC & total & 0.0014 (1/1400)  & 0.00071 (0/1400) \\
    GENIE~2.12.10 + MEC & shape & 0.0014 (1/1400)  & 0.0021 (2/1400) \\
    NuWro~21.09         & total & 0.00056 (0/1800) & 0.00056 (0/1800) \\
    NuWro~21.09         & shape & 0.00056 (0/1800) & 0.0033 (5/1800) \\
    GiBUU~2019          & total & 0.00071 (0/1400) & 0.00071 (0/1400) \\
    GiBUU~2019          & shape & 0.00071 (0/1400) & 0.00071 (0/1400) \\
    \bottomrule
  \end{tabular}
\end{table}

\paragraph{Hadronic-response condition and recoil-energy sensitivity.}
% [records reorg] former heading: \subsubsection{Hadronic-response condition and recoil-energy sensitivity}
MINERvA's recoil-energy analyses assign a hadronic-energy response uncertainty derived from test-beam calorimetry~\cite{MINERvA:2015testbeam}, separate from the GEANT hadron-interaction reweights~\cite{MINERvA:2021owq,MINERvA:2022incl}. In the low-recoil measurement the resulting hadronic-energy uncertainty rises to \SI{10}{\percent} at its highest three-momentum transfers, $0.9<q_{3}<\SI{1.2}{GeV}$~\cite{MINERvA:2022incl}. The specified nuisance model here includes the reweights but no calorimetric response band. Such a band would act on all five nulls together, through the shared detector simulation and unfolding. A response mismodeling could move each observed statistic in either direction, and a band drawn in the calibration would also broaden and could shift each null distribution. The net effect on each $p$ value is unmeasured. The GENIE~2.12.10 projection result of \S\ref{sec:joint5d-projections} shows that the information rejecting that prediction involves the hadronic axes, on which a response band would act.
% Search: net effect on each; calorimetric response band; response band drawn @ 1a015eef
In an exploratory sensitivity study of the response~\cite{JointTestRecoil2026}, the data are re-unfolded with the simulation's reconstructed recoil energy scaled by $\pm\SI{4}{\percent}$, and the ten tests are re-evaluated against the null distributions used for the decisions. This scale is the level of data--simulation agreement reported for test-beam protons, pions and electrons of 0.35--\SI{2.0}{GeV/c}~\cite{MINERvA:2015testbeam}. Neutrons and higher-momentum hadrons, which matter at high $W$, lie outside that sample. The $\pm4\%$ scale is not a validated uncertainty prescription. At both $-4\%$ and $+4\%$ every one of the ten rejections holds. Nine tests keep $k=0$. NuWro's shape test, the closest to its null, changes from $k=1$ to $k=3$ of 1751 at $+4\%$ ($p=0.0023$, against a Holm threshold of 0.05). The null distributions were not regenerated with the variation, although they are simulated through the same detector model, so the size of this omission is unmeasured for these strongly non-central nulls. The study also leaves untested particle-species, energy-dependent or resolution effects, a response difference applied before calibration, and larger variations. The rejections are therefore conditional on the simulation's hadronic response, and the response model is not shown to be complete. The observed statistics at each scale, and the study's controls, are given in the records document.
% absence-ok: "not regenerated" and "leaves untested" scope what this exploratory study varied (condensed from the paragraph moved verbatim to the records, JE-12); they are not a claim that a result does not exist.

\paragraph{Recoil-energy scale study.}
% [records reorg, repeated sentence not printed] In an exploratory sensitivity study of the response~\cite{JointTestRecoil2026}, the data are re-unfolded with the simulation's reconstructed recoil energy scaled by $\pm\SI{4}{\percent}$, and the ten tests are re-evaluated against the null distributions used for the decisions.
% [records reorg, repeated sentence not printed] This scale is the level of data--simulation agreement reported for test-beam protons, pions and electrons of 0.35--\SI{2.0}{GeV/c}~\cite{MINERvA:2015testbeam}.
% [records reorg, repeated sentence not printed] Neutrons and higher-momentum hadrons, which matter at high $W$, lie outside that sample.
% [records reorg, repeated sentence not printed] The $\pm4\%$ scale is not a validated uncertainty prescription.
% [records reorg, repeated sentence not printed] At both $-4\%$ and $+4\%$ every one of the ten rejections holds.
% [records reorg, repeated sentence not printed] Nine tests keep $k=0$.
% [records reorg, repeated sentence not printed] NuWro's shape test, the closest to its null, changes from $k=1$ to $k=3$ of 1751 at $+4\%$ ($p=0.0023$, against a Holm threshold of 0.05).
 A re-unfold with unscaled recoil reproduces the decisions of the earlier note's table of the joint generator tests, with every claim $p$ value inside the range spanned by the data's rounding variations. The observed statistics move by up to about \SI{11}{\percent} from their recorded values. The evaluation was reproduced independently. Table~\ref{tab:joint5d-recoil} gives the observed statistics.
% Search: validated uncertainty prescription; calorimetric response band; recoil-energy-scale band @ ddf16f10
% Deposit pass: W2B-REPORT-20261006.md sec. 5, "Observed statistics" (descriptive) and the k/B table; values copied as recorded. Null medians are the frozen unshifted (c = 0) null medians.
\begin{table}[htbp]
  \caption[Descriptive recoil-scale sensitivity of the observed statistics]{All ten tests remain rejected when the simulation's reconstructed recoil energy is scaled by $\pm4\%$ against unchanged null distributions. The observed joint statistics are shown beside the recorded values and the median of each unshifted null distribution, which is not regenerated. Every test keeps $k=0$ at both signs except NuWro's shape test ($k=1$ of 1751 at $-4\%$, 3 at $+4\%$). The values are descriptive.}
  \label{tab:joint5d-recoil}
  \centering\small
  \begin{tabular}{@{}lrrrrrrrr@{}}
    \toprule
    & \multicolumn{4}{c}{$T_{\mathrm{tot}}$} & \multicolumn{4}{c}{$T_{\mathrm{sh}}$} \\
    \cmidrule(lr){2-5}\cmidrule(lr){6-9}
    $G$ & recorded & $-4\%$ & $+4\%$ & null median & recorded & $-4\%$ & $+4\%$ & null median \\
    \midrule
    MINERvA Tune~v1     & 2852 & 2856 & 2540 & 49   & 631  & 647  & 600  & 49 \\
    GENIE~2.12.10       & 1813 & 1796 & 1721 & 528  & 1086 & 1104 & 1040 & 582 \\
    GENIE~2.12.10 + MEC & 1647 & 1652 & 1595 & 406  & 869  & 893  & 839  & 450 \\
    NuWro~21.09         & 7427 & 7380 & 7361 & 6324 & 9660 & 9690 & 9540 & 7110 \\
    GiBUU~2019          & 2867 & 3129 & 2665 & 1791 & 4013 & 4154 & 3901 & 1975 \\
    \bottomrule
  \end{tabular}
\end{table}
% [records reorg, repeated sentence not printed] The rejections are therefore conditional on the simulation's hadronic response, and the response model is not shown to be complete.
% [records reorg, repeated sentence not printed] It leaves untested (a) the response of the null distributions, which were not regenerated with the variation although they are simulated through the same detector model, so the size of this omission is unmeasured for these strongly non-central nulls;
The study measures the sensitivity of the observed statistics and decisions to one coherent scale of reconstructed available energy and recoil energy, applied after the energy calibration at a single magnitude. (b) particle-species, energy-dependent or resolution effects, and a response difference applied before calibration; and (c) larger variations.
% Register round n1 (2026-10-25): this subsection is the one home of the recoil-response omission; the condition
% (ii) gloss in sec:joint5d-conditions points here. The fixed null distributions print in the paragraph above and
% in (a); the need to vary them in a response-uncertainty construction, in the first paragraph.
% Search: regenerated with the; response-uncertainty; calorimetric response band; net effect on each p value @ ddf16f10

\paragraph{The published comparison in muon kinematics.}
The published comparison in muon kinematics alone, whose observables do not use the reconstructed recoil energy, also disfavors the same generator families~\cite{MINERvA:2021owq}. It provides context for a disagreement involving muon observables, but the hypotheses tested here use a different estimator and calibration.

\paragraph{Hydrocarbon target correction.}
NuWro~21.09 and GiBUU~2019 were generated on carbon, while the data and the other predictions are per nucleon of CH. This study corrects the two predictions to CH and re-evaluates their tests against the null distributions used for the decisions. These target-corrected evaluations are descriptive. The decisions of the earlier note's table of the joint generator tests are unchanged. The primary generator comparisons (\S\ref{sec:fived16-gen}) do not use this correction and show NuWro and GiBUU on carbon only. This section records the correction.
% Descriptive, decisions unchanged: DECISION-20261020-joseph-nuwro-ch-claim.md:37 ("the replays are descriptive and
% do not re-grade A7"), reworded on the director's ruling of 2026-10-25 (register round n1).

\paragraph{Construction.}
% [records reorg] former heading: \paragraph{Construction.}
 Each prediction is rebuilt per fine cell as $x_{\mathrm{CH}}=\tfrac{12}{13}\,x_{\mathrm{C}}+\tfrac{1}{13}\,x_{\mathrm{H}}$. Here $x_{\mathrm{H}}$ is the generator's own free-proton product, run through the unchanged conversion, flux-repair and binning chain. Each prediction's Monte Carlo variance is combined with the squares of the same weights. The hydrogen fraction is $f_{\mathrm{H}}=5.08\%$ for NuWro over the grid and 4.77\% for GiBUU over its 72 cells. Over the reported cells, the ratio $r_{c}=\mu_{\mathrm{CH}}/\mu_{\mathrm{C}}$ has minimum, median, and maximum 0.923, 0.963, and 1.039 for NuWro (109 cells) and 0.923, 0.963, and 1.036 for GiBUU (72 cells). Per cell, $f_{\mathrm{H}}$ reaches 0.112 and 0.109.  On these cells both corrections move the prediction toward the data. Its effect on the sixteen-cell comparison is stated in \S\ref{sec:fived16-gen-target}.
% B:29-34 (combination, sumw2, the frozen chain, the supplement); C:46-51 (GiBUU, frozen r(E), 72 cells).
% Accuracy fix (unscored; director 2026-10-26): "with the same weights" misstated the method. The code fills each
% event with weight c = w r / N x (12/13 or 1/13) and sums c^2 for the variance (docs/orchestration/state/
% measurement-paper/phase1/mp2/s2-receipts/d4/harmonize_9dbcde93.py:71, :104-108, :312, :418, :426 @ 7ad668f4), so
% sumw2_CH = (12/13)^2 sumw2_C12 + (1/13)^2 sumw2_H (mp3b/REPORT.md:30).
% Draft 6 (REVIEW-2 finding 12): B:97 (grid integral; 0.0506 over the 109 cells, B:105), C:51 (72-cell domain).
% B:97, :103-105; C:51, :56-58.
% B:138-139; C:80-81 (the cross term is negative: the data lie on the CH side). REVIEW-2 finding 11: pointer moved here.

\paragraph{Re-evaluated tests.}
% [records reorg] former heading: \paragraph{Re-evaluated tests.}
 With the metric used for the decisions, NuWro's total statistic falls by 3.83 null standard deviations and its shape statistic by 0.37.  At its original Holm position (threshold 0.005) the total test would then lie 1.11 null standard deviations below its boundary, and so be undetermined.  It remains rejected only because its larger $p$ value moves it to the last Holm position (threshold 0.05). In the $\kappa=\jtMoneKappaRobust$ family, which is for reporting only, it is not rejected ($p=0.075$). NuWro's shape test stays rejected, 0.74 null standard deviations above its decision boundary.  GiBUU's statistics fall by 4.67 (total) and 1.33 (shape) null standard deviations. Both tests stay rejected at their original Holm positions, including at $\kappa=\jtMoneKappaRobust$. Their margins shrink from 7.86 to 3.18 and from 3.58 to 2.24 null standard deviations.  With each CH prediction's own variance in the metric, no decision changes for either generator.
% B:123-137 (NuWro: T_total 7427.1 -> 6759.8, -3.83 null SD; T_shape 9660.0 -> 9333.1, -0.37; frozen position:
% threshold 0.005, T_crit 6953.94, -1.11 SD, undetermined; re-sorted to the last position, threshold 0.05, CP
% [0.0183, 0.0336], rejected; shape to threshold 0.025, 0.74 SD above its frozen T_crit; kappa = 3: total p 0.0748,
% k = 130, not rejected; shape rejected). C:21, :72-79 (GiBUU: T_total -4.67, T_shape -1.33 null SD; margins 7.86 ->
% 3.18 and 3.58 -> 2.24; frozen positions, no re-sort; kappa = 3 also rejected). M:65-78 (frozen margins).
% B:129 (W_CH: both rejected, no change); C:76.

\paragraph{Generation checks.}
% [records reorg] former heading: \paragraph{Generation checks.}
 For GiBUU, the sampling-closure check failed on the hydrogen sample. A bias check of that check's premise was added post hoc. Its criterion, a change of at most $0.1\,\sigma_{c}$ in every cell, was specified before the check was computed. The check met the criterion: it moved the CH prediction by at most $5.1\times10^{-4}\,\sigma_{c}$. Here $\sigma_{c}$ is the standard deviation of the cell's result over NuWro's null distribution used for the decisions ($B=1751$). The GiBUU re-evaluation is reported under the label ``B3 (sampling closure) failed on hydrogen; a post-hoc bias check of its premise passed at $\le0.1\,\sigma_{c}$''.
It has a further limitation: GiBUU's deep-inelastic events on hydrogen, 46\% of its hydrogen sample, do not conserve energy.  They lose a median of \SI{45}{MeV}, with a median transverse-momentum imbalance of \SI{31}{MeV}, while its other hydrogen events balance exactly. The cause has not been identified. The available energy of those events may be low, and the non-conservation may also affect GiBUU's carbon prediction, which cannot be checked in this way.
% [records reorg] carried verbatim from the duplicate 2cea6ebd64-2: the only sentence of the dropped copy of "Generation checks" that the kept copy lacks (D2)
NuWro's 50--\SI{100}{GeV} hydrogen supplement is flagged rather than assumed negligible.
In NuWro's hydrogen generation, four chunks aborted and were rerun on alternate seeds. The fraction of events this excludes is estimated as $1.8\times10^{-5}$ in the main sample and $3.8\times10^{-5}$ in the 50--\SI{100}{GeV} supplement, or $1.9\times10^{-5}$ and $6.2\times10^{-5}$ including the first attempt. These estimates rest on four to six events.  The supplement's 95\% upper bound of $1.8\times10^{-4}$ exceeds the $10^{-4}$ level set for it, so the supplement is flagged rather than assumed negligible.
% C:7-8 (the required label, verbatim), :23-44 (declaration 00efbb25..., recorded before computing; max |Delta_c/sigma_c|
% 5.1e-4 against the 0.1 acceptance). B:48-56 (NuWro hydrogen chunks: 3 + 1 aborts rerun on alternate seeds; excluded
% fraction 1.8e-5 main, 3.8e-5 supplement; the supplement's 95% bound 1.8e-4 above the 1e-4 bar, FLAGGED).
% sigma_c: mp3c/REPORT.md:31 takes sigma_c from MP3-B; mp3b/REPORT.md:37, reading (a), "the SD of the J-cell result
% over NuWro's frozen B = 1751 null"; max |Delta_c/sigma_c| 5.1e-4 in reading (a), mp3c/REPORT.md:38.
% Draft 6: docs/orchestration/state/measurement-paper/BINDING-PHASE1-ADDENDUM-9.md:28-31 (sec. 4, "Recorded for the record";
% in tree since the rebase onto 51da2a62): production_ID 34, 46% of the GiBUU-H sample; median energy loss
% 45 MeV, median p_T imbalance 31 MeV; non-DIS hydrogen events balance exactly; cause not identified.
% ADDENDUM-9 :29-31 ("Carbon cannot be checked this way"; "possibly on the frozen GiBUU product").
% Search: do not conserve energy; energy non-conservation; production_ID 34 (cause not identified) @ ddf16f10

% [records reorg] former heading: \subsubsection{Hydrocarbon target correction for NuWro and GiBUU}
% [records reorg, duplicate not printed] NuWro~21.09 and GiBUU~2019 were generated on carbon, while the data and the other predictions are per nucleon of CH. This study corrects the two predictions to CH, per fine cell as $x_{\mathrm{CH}}=\tfrac{12}{13}\,x_{\mathrm{C}}+\tfrac{1}{13}\,x_{\mathrm{H}}$ with $x_{\mathrm{H}}$ the generator's own free-proton product, and re-evaluates their tests against the null distributions used for the decisions. These target-corrected evaluations are descriptive. The decisions of the earlier note's table of the joint generator tests are unchanged. The primary generator comparisons (\S\ref{sec:fived16-gen}) do not use this correction and show NuWro and GiBUU on carbon only. On the reported cells both corrections move the prediction toward the data. Its effect on the sixteen-cell comparison is stated in \S\ref{sec:fived16-gen-target}.
% [records reorg, duplicate not printed] With the metric used for the decisions, NuWro's total statistic falls by 3.83 null standard deviations. At its original Holm position (threshold 0.005) the total test would then lie 1.11 null standard deviations below its boundary, and so be undetermined. It remains rejected only because its larger $p$ value moves it to the last Holm position (threshold 0.05). In the $\kappa=\jtMoneKappaRobust$ family, which is for reporting only, it is not rejected ($p=0.075$). NuWro's shape test stays rejected. Both GiBUU tests stay rejected at their original Holm positions, including at $\kappa=\jtMoneKappaRobust$. With each CH prediction's own variance in the metric, no decision changes for either generator.

% [records reorg] former heading: \paragraph{Generation checks.}
% [records reorg, duplicate not printed]  For GiBUU, the sampling-closure check failed on the hydrogen sample. A bias check of that check's premise was added post hoc. Its criterion, a change of at most $0.1\,\sigma_{c}$ in every cell, was specified before the check was computed. The check met the criterion: it moved the CH prediction by at most $5.1\times10^{-4}\,\sigma_{c}$. Here $\sigma_{c}$ is the standard deviation of the cell's result over NuWro's null distribution used for the decisions ($B=1751$). The GiBUU re-evaluation is reported under the label ``B3 (sampling closure) failed on hydrogen; a post-hoc bias check of its premise passed at $\le0.1\,\sigma_{c}$''.
% [records reorg, duplicate not printed] The correction carries two further generation caveats. GiBUU's deep-inelastic events on hydrogen, 46\% of its hydrogen sample, do not conserve energy. The cause has not been identified, and the non-conservation may also affect GiBUU's carbon prediction, which cannot be checked in this way. NuWro's 50--\SI{100}{GeV} hydrogen supplement is flagged rather than assumed negligible. The construction's details, the margins and the numbers of the generation checks are given in the records document.
